% !TEX program = lualatex
% !TEX encoding = UTF-8 Unicode
%%%%%%%%%%%%%%%%%%%%%%%%%%%%%%%%%%%%%%%%%%%%%%%%%%%%%%%%%
% IWES Abschlussarbeit Vorlage (Deutsch)
% Version:  2.0.0
% Datum:    2026-03-09
% Autor:    Basem Rajjoub
% Email:    basem.rajjoub@iwes.uni-hannover.de
%%%%%%%%%%%%%%%%%%%%%%%%%%%%%%%%%%%%%%%%%%%%%%%%%%%%%%%%%
%
% Hauptdokument - Diese Datei kompilieren.
%
% COMPILER:
%   lualatex 00_main && biber 00_main && lualatex 00_main
%%%%%%%%%%%%%%%%%%%%%%%%%%%%%%%%%%%%%%%%%%%%%%%%%%%%%%%%%


\documentclass[
11pt,                   % Schriftgröße 11pt
a4paper,                % Papierformat A4
twoside,                % Zweiseitiges Dokument (für 'oneside' ändern)
openright,              % Kapitel beginnen auf rechten (ungeraden) Seiten
cleardoublepage=empty,  % Leerseiten ohne Kopfzeilen
parskip=half,           % Absatzabstand halbe Zeile
numbers=noenddot,       % Keine Punkte am Ende der Nummerierung
listof=totoc,           % Verzeichnisse im Inhaltsverzeichnis
bibliography=totoc,     % Literaturverzeichnis im Inhaltsverzeichnis
]{scrreprt}


% Einstellungen und Pakete laden
% Einstellungen und Pakete für die IWES Vorlage
% Unterstützt: LuaLaTeX (Rotis-Schrift) und pdfLaTeX (Cabin-Fallback)

%=== Kompatibilität ====================================
\usepackage{scrhack}    % KOMA-Script Kompatibilität

%=== Schriften =========================================
\usepackage{iftex}

\ifluatex
    \usepackage{fontspec}
    \IfFileExists{./font/RotisSansSerifStd-Regular.otf}{%
        \setsansfont{RotisSansSerifStd-Regular.otf}[
            Path           = ./font/,
            BoldFont       = RotisSansSerifStd-Bold.otf,
            ItalicFont     = RotisSansSerifStd-Italic.otf,
            BoldItalicFont = RotisSansSerifStd-Bold.otf,
            SmallCapsFont  = RotisSansSerifStd-Regular.otf,
            SmallCapsFeatures = {Scale=0.8},
        ]
        \renewcommand*\familydefault{\sfdefault}
    }{%
        \setsansfont{TeX Gyre Heros}
        \renewcommand*\familydefault{\sfdefault}
        \PackageWarning{IWES}{Rotis nicht gefunden - verwende TeX Gyre Heros}
    }
\else
    \usepackage[T1]{fontenc}
    \usepackage[utf8]{inputenc}
    \usepackage[sfdefault]{cabin}
\fi

%=== Sprache ===========================================
\usepackage[english, ngerman]{babel}    % Beide geladen; aktive Sprache wird per \selectlanguage in 00_main.tex gesetzt
\usepackage{csquotes}       % Für biblatex erforderlich
\usepackage{etoolbox}       % Für Bedingungs-Makros (\ifdefstring)
\usepackage{hyphenat}

%=== Arbeitstyp-Konfiguration ===========================
% Befehl: \setthesistype{master|bachelor|ip|wai}
% Setzt automatisch alle typ-spezifischen Variablen
\newcommand{\setthesistype}[1]{%
    \def\thesistype{#1}%
    % Master oder Bachelor
    \ifboolexpr{ test {\ifdefstring{\thesistype}{master}} or test {\ifdefstring{\thesistype}{bachelor}} }{%
        \def\thesisdesc{Wissenschaftliche Arbeit zur Erlangung des akademischen Grades}%
        \def\thesismodule{}%
        \newbool{showcoexaminer}\booltrue{showcoexaminer}%
        \newbool{showdegree}\booltrue{showdegree}%
    }{}%
    % Interdisziplinäres Projekt
    \ifdefstring{\thesistype}{ip}{%
        \def\thesisdesc{Ausarbeitung für das Modul}%
        \def\thesismodule{Interdisziplinäres Projekt}%
        \newbool{showcoexaminer}\boolfalse{showcoexaminer}%
        \newbool{showdegree}\boolfalse{showdegree}%
    }{}%
    % Wissenschaftliches Arbeiten im Ingenieurwesen
    \ifdefstring{\thesistype}{wai}{%
        \def\thesisdesc{Ausarbeitung für das Modul}%
        \def\thesismodule{Wissenschaftliches Arbeiten im Ingenieurwesen}%
        \newbool{showcoexaminer}\boolfalse{showcoexaminer}%
        \newbool{showdegree}\boolfalse{showdegree}%
    }{}%
}

%=== Seitenlayout ======================================
\usepackage[
    a4paper,
    includeheadfoot,
    left=2.5cm,
    right=2.5cm,
    top=2cm,
    bottom=0.5cm
]{geometry}

\usepackage{pdfpages}       % PDF-Dateien einbinden

%=== Kopf-/Fußzeilen ===================================
\usepackage[headsepline, automark]{scrlayer-scrpage}
\clearpairofpagestyles
\automark[section]{chapter}

\makeatletter
\if@twoside
    \lehead{\pagemark}
    \rehead{\slshape\leftmark}
    \lohead{\slshape\rightmark}
    \rohead{\pagemark}
\else
    \lohead{\slshape\headmark}
    \rohead{\pagemark}
\fi
\makeatother

\pagestyle{scrheadings}

% Kompatibilitätsmakros: \lhead{}/\rhead{} in 00_main.tex sind Kurzbefehle für
% laufende Kolumnentitel. Es handelt sich NICHT um fancyhdr-Befehle — sie rufen
% intern \markboth/\markright von KOMA-Script auf.
\DeclareDocumentCommand{\lhead}{m}{\markboth{#1}{#1}}
\DeclareDocumentCommand{\rhead}{m}{\markright{#1}}
\DeclareDocumentCommand{\chead}{m}{}

%=== Institutionelle Konstanten ========================
% Diese Werte anpassen, falls sich Institutsname, Professor oder Adresse ändern.
\def\universityName{Leibniz Universität Hannover}
\def\instituteName{Institut für Windenergiesysteme}
\def\professorName{Prof. Dr.-Ing. Andreas Reuter}
\def\instituteStreet{Appelstraße 9A}
\def\instituteCity{30167 Hannover}

%=== Zeilenabstand =====================================
\usepackage{setspace}
\onehalfspacing

%=== Hyperlinks ========================================
\ifluatex
    \usepackage[verbose=silent]{microtype}
\else
    \usepackage{microtype}
\fi

\usepackage{hyperref}
\hypersetup{
    hypertexnames=false,
    colorlinks=true,
    linkcolor=blue!70!black,
    citecolor=green!50!black,
    urlcolor=blue,
    filecolor=magenta!70!black,
}

\usepackage{xurl}           % Besserer URL-Umbruch

%=== Tabellen ==========================================
\usepackage{colortbl}
\usepackage{booktabs}       % Professionelle Tabellen
\usepackage{tabularx}
\usepackage{longtable}
\usepackage{array}
\newcolumntype{P}[1]{>{\centering\arraybackslash}p{#1}}
\usepackage{rotating}
\usepackage{multirow}
\usepackage{multicol}

%=== KI-Dokumentation ==================================
% Spaltentyp und Boxen für die Dokumentation des KI-Einsatzes (91_ki_dokumentation.tex)
\newcolumntype{L}[1]{>{\raggedright\arraybackslash}p{#1}}
\usepackage{tcolorbox}
\tcbuselibrary{breakable, skins}
\definecolor{kiakzent}{rgb}{0.0,0.314,0.608}   % LUH-Blau
% Prompt-Box:   \begin{kiprompt}[Titel] ... \end{kiprompt}
\newtcolorbox{kiprompt}[1][Prompt]{%
    enhanced, breakable, sharp corners,
    colback=kiakzent!6, colframe=kiakzent, boxrule=0pt, leftrule=2.5pt,
    left=6pt, right=6pt, top=4pt, bottom=4pt,
    fonttitle=\bfseries\footnotesize, coltitle=kiakzent, colbacktitle=kiakzent!6,
    title={#1}, attach title to upper={\par\smallskip}, fontupper=\small,
    before skip=6pt, after skip=4pt,
}
% Antwort-Box:  \begin{kiantwort}[Titel] ... \end{kiantwort}
\newtcolorbox{kiantwort}[1][Antwort]{%
    enhanced, breakable, sharp corners,
    colback=black!3, colframe=black!45, boxrule=0pt, leftrule=2.5pt,
    left=6pt, right=6pt, top=4pt, bottom=4pt,
    fonttitle=\bfseries\footnotesize, coltitle=black!60,
    title={#1}, attach title to upper={\par\smallskip}, fontupper=\small,
    before skip=4pt, after skip=10pt,
}

%=== Grafiken ==========================================
\usepackage{graphicx}
\graphicspath{{figures/}}   % Bilderpfad
\usepackage{float}
\usepackage{wrapfig}
\usepackage{caption}
\usepackage{subcaption}
\usepackage[section]{placeins}

%=== Mathematik ========================================
\usepackage{amsmath}
\usepackage{amsfonts}
\usepackage{amssymb}

%=== Listen ============================================
\usepackage{enumitem}       % Flexible Listenabstände (ersetzt paralist)
\setlist{nosep}             % Kompakte Listen global (kein zusätzlicher vertikaler Abstand)

%=== Einheiten und Symbole =============================
\usepackage{siunitx}
\sisetup{locale=DE}         % Deutsche Formatierung
\usepackage[right]{eurosym}

%=== Typographie =======================================
\clubpenalty = 10000
\widowpenalty = 10000
\displaywidowpenalty = 10000
\tolerance = 10000
\hbadness = 10000
\emergencystretch = 3em
\hfuzz = 1pt

%=== Inhaltsverzeichnis ================================
\RedeclareSectionCommand[tocpagenumberwidth=2.5em]{chapter}

%=== Literaturverzeichnis (biblatex) ===================
\usepackage[
    backend=biber,
    style=numeric,
    sorting=none,
    doi=true,
    url=true,
    isbn=true,
    maxbibnames=99,
    giveninits=true,
]{biblatex}

\addbibresource{99_bibliography.bib}

% DOI als klickbaren Link formatieren
\DeclareFieldFormat{doi}{%
    \mkbibacro{DOI}\addcolon\space
    \ifhyperref
        {\href{https://doi.org/#1}{\nolinkurl{#1}}}
        {\nolinkurl{#1}}%
}

% Warnungen bei fehlenden DOI/URL/ISBN
\AtEveryBibitem{%
    % Artikel: DOI oder URL erforderlich
    \ifentrytype{article}{%
        \iffieldundef{doi}{%
            \iffieldundef{url}{%
                \PackageWarning{biblatex}{Fehlende DOI/URL für Artikel: \thefield{entrykey}}%
            }{}%
        }{}%
    }{}%
    % Bücher: ISBN oder DOI erforderlich
    \ifentrytype{book}{%
        \iffieldundef{isbn}{%
            \iffieldundef{doi}{%
                \PackageWarning{biblatex}{Fehlende ISBN/DOI für Buch: \thefield{entrykey}}%
            }{}%
        }{}%
    }{}%
}



%%%%%%%%%%%%%%%%%%%%%%%%%%%%%%%%%%%%%%%%%%%%%%%%%%%%%%%%%
%   STUDIERENDEN-KONFIGURATION — NUR DIESEN BEREICH BEARBEITEN
%   Alles unterhalb von \begin{document} muss normalerweise nicht geändert werden
%   (außer Kapitel-\input-Zeilen hinzufügen/entfernen)
%%%%%%%%%%%%%%%%%%%%%%%%%%%%%%%%%%%%%%%%%%%%%%%%%%%%%%%%%

% PFLICHTFELD: Art der Arbeit
%%%%%%%%%%%%%%%%%%%%%%%%%%%%%%%%%%%%%%%%%%%%%%%%%%%%%%%%%
% Gueltige Werte: master, bachelor, ip, wai
%   master   = Masterarbeit
%   bachelor = Bachelorarbeit
%   ip       = Interdisziplinaeres Projekt
%   wai      = Wissenschaftliches Arbeiten im Ingenieurwesen
%%%%%%%%%%%%%%%%%%%%%%%%%%%%%%%%%%%%%%%%%%%%%%%%%%%%%%%%

\setthesistype{master}

%%%%%%%%%%%%%%%%%%%%%%%%%%%%%%%%%%%%%%%%%%%%%%%%%%%%%%%%
% Persoenliche Angaben fuer das Titelblatt
%%%%%%%%%%%%%%%%%%%%%%%%%%%%%%%%%%%%%%%%%%%%%%%%%%%%%%%%

\def\name{Max Mustermann}                           % Name des Autors
\def\titel{Beispieltitel}        % Titel der Arbeit
\def\matrikelnummer{12345678}                       % Matrikelnummer

\def\projectType{Master of Science (M.Sc.)}         % Abschlussgrad (fuer master/bachelor)
% \def\projectType{Bachelor of Science (B.Sc.)}

\def\betreuer{Titel Vorname Nachname}               % Betreuer
\def\fbetreuer{Institut für Windenergiesysteme, LUH}

\def\pruefer{Titel Vorname Nachname}                % Erstprüfer
\def\fpruefer{Institut für Windenergiesysteme, LUH}

\def\zweitpruefer{Titel Vorname Nachname}           % Zweitprüfer
\def\fzweitpruefer{Institut für Windenergiesysteme, LUH}

\def\studiengang{Bauingenieurwesen (M.Sc.)}         % Studiengang
\def\projectNumber{XXX-XXX-XXX}                     % Projektnummer

% Nur bei vertraulichen Arbeiten: Firmenname für den Sperrvermerk setzen.
% Wenn die Arbeit nicht vertraulich ist, den gesamten Sperrvermerk-Abschnitt in 03_erklaerungen.tex auskommentieren.
\def\companyname{Firmaxxx}


%%%%%%%%%%%%%%%%%%%%%%%%%%%%%%%%%%%%%%%%%%%%%%%%%%%%%%%%
\begin{document}
%%%%%%%%%%%%%%%%%%%%%%%%%%%%%%%%%%%%%%%%%%%%%%%%%%%%%%%%

\selectlanguage{ngerman}

%=== Titelseiten ========================================
% Automatisch ausgewaehlt basierend auf \thesistype
% ============================================================================
% Titelseiten
% ============================================================================
%
% WICHTIG: \setthesistype{...} muss in 00_main.tex aufgerufen werden!
% Gueltige Werte: master, bachelor, ip, wai
%
% ============================================================================

% Pruefe ob \thesistype definiert ist
\ifdef{\thesistype}{}{%
    \PackageError{IWES}{%
        Arbeitstyp nicht definiert!^^J%
        Bitte in 00_main.tex setzen: \string\setthesistype{master}^^J%
        Gueltige Werte: master, bachelor, ip, wai%
    }{Setze \string\setthesistype in 00_main.tex nach \string% Einstellungen und Pakete für die IWES Vorlage
% Unterstützt: LuaLaTeX (Rotis-Schrift) und pdfLaTeX (Cabin-Fallback)

%=== Kompatibilität ====================================
\usepackage{scrhack}    % KOMA-Script Kompatibilität

%=== Schriften =========================================
\usepackage{iftex}

\ifluatex
    \usepackage{fontspec}
    \IfFileExists{./font/RotisSansSerifStd-Regular.otf}{%
        \setsansfont{RotisSansSerifStd-Regular.otf}[
            Path           = ./font/,
            BoldFont       = RotisSansSerifStd-Bold.otf,
            ItalicFont     = RotisSansSerifStd-Italic.otf,
            BoldItalicFont = RotisSansSerifStd-Bold.otf,
            SmallCapsFont  = RotisSansSerifStd-Regular.otf,
            SmallCapsFeatures = {Scale=0.8},
        ]
        \renewcommand*\familydefault{\sfdefault}
    }{%
        \setsansfont{TeX Gyre Heros}
        \renewcommand*\familydefault{\sfdefault}
        \PackageWarning{IWES}{Rotis nicht gefunden - verwende TeX Gyre Heros}
    }
\else
    \usepackage[T1]{fontenc}
    \usepackage[utf8]{inputenc}
    \usepackage[sfdefault]{cabin}
\fi

%=== Sprache ===========================================
\usepackage[english, ngerman]{babel}    % Beide geladen; aktive Sprache wird per \selectlanguage in 00_main.tex gesetzt
\usepackage{csquotes}       % Für biblatex erforderlich
\usepackage{etoolbox}       % Für Bedingungs-Makros (\ifdefstring)
\usepackage{hyphenat}

%=== Arbeitstyp-Konfiguration ===========================
% Befehl: \setthesistype{master|bachelor|ip|wai}
% Setzt automatisch alle typ-spezifischen Variablen
\newcommand{\setthesistype}[1]{%
    \def\thesistype{#1}%
    % Master oder Bachelor
    \ifboolexpr{ test {\ifdefstring{\thesistype}{master}} or test {\ifdefstring{\thesistype}{bachelor}} }{%
        \def\thesisdesc{Wissenschaftliche Arbeit zur Erlangung des akademischen Grades}%
        \def\thesismodule{}%
        \newbool{showcoexaminer}\booltrue{showcoexaminer}%
        \newbool{showdegree}\booltrue{showdegree}%
    }{}%
    % Interdisziplinäres Projekt
    \ifdefstring{\thesistype}{ip}{%
        \def\thesisdesc{Ausarbeitung für das Modul}%
        \def\thesismodule{Interdisziplinäres Projekt}%
        \newbool{showcoexaminer}\boolfalse{showcoexaminer}%
        \newbool{showdegree}\boolfalse{showdegree}%
    }{}%
    % Wissenschaftliches Arbeiten im Ingenieurwesen
    \ifdefstring{\thesistype}{wai}{%
        \def\thesisdesc{Ausarbeitung für das Modul}%
        \def\thesismodule{Wissenschaftliches Arbeiten im Ingenieurwesen}%
        \newbool{showcoexaminer}\boolfalse{showcoexaminer}%
        \newbool{showdegree}\boolfalse{showdegree}%
    }{}%
}

%=== Seitenlayout ======================================
\usepackage[
    a4paper,
    includeheadfoot,
    left=2.5cm,
    right=2.5cm,
    top=2cm,
    bottom=0.5cm
]{geometry}

\usepackage{pdfpages}       % PDF-Dateien einbinden

%=== Kopf-/Fußzeilen ===================================
\usepackage[headsepline, automark]{scrlayer-scrpage}
\clearpairofpagestyles
\automark[section]{chapter}

\makeatletter
\if@twoside
    \lehead{\pagemark}
    \rehead{\slshape\leftmark}
    \lohead{\slshape\rightmark}
    \rohead{\pagemark}
\else
    \lohead{\slshape\headmark}
    \rohead{\pagemark}
\fi
\makeatother

\pagestyle{scrheadings}

% Kompatibilitätsmakros: \lhead{}/\rhead{} in 00_main.tex sind Kurzbefehle für
% laufende Kolumnentitel. Es handelt sich NICHT um fancyhdr-Befehle — sie rufen
% intern \markboth/\markright von KOMA-Script auf.
\DeclareDocumentCommand{\lhead}{m}{\markboth{#1}{#1}}
\DeclareDocumentCommand{\rhead}{m}{\markright{#1}}
\DeclareDocumentCommand{\chead}{m}{}

%=== Institutionelle Konstanten ========================
% Diese Werte anpassen, falls sich Institutsname, Professor oder Adresse ändern.
\def\universityName{Leibniz Universität Hannover}
\def\instituteName{Institut für Windenergiesysteme}
\def\professorName{Prof. Dr.-Ing. Andreas Reuter}
\def\instituteStreet{Appelstraße 9A}
\def\instituteCity{30167 Hannover}

%=== Zeilenabstand =====================================
\usepackage{setspace}
\onehalfspacing

%=== Hyperlinks ========================================
\ifluatex
    \usepackage[verbose=silent]{microtype}
\else
    \usepackage{microtype}
\fi

\usepackage{hyperref}
\hypersetup{
    hypertexnames=false,
    colorlinks=true,
    linkcolor=blue!70!black,
    citecolor=green!50!black,
    urlcolor=blue,
    filecolor=magenta!70!black,
}

\usepackage{xurl}           % Besserer URL-Umbruch

%=== Tabellen ==========================================
\usepackage{colortbl}
\usepackage{booktabs}       % Professionelle Tabellen
\usepackage{tabularx}
\usepackage{longtable}
\usepackage{array}
\newcolumntype{P}[1]{>{\centering\arraybackslash}p{#1}}
\usepackage{rotating}
\usepackage{multirow}
\usepackage{multicol}

%=== KI-Dokumentation ==================================
% Spaltentyp und Boxen für die Dokumentation des KI-Einsatzes (91_ki_dokumentation.tex)
\newcolumntype{L}[1]{>{\raggedright\arraybackslash}p{#1}}
\usepackage{tcolorbox}
\tcbuselibrary{breakable, skins}
\definecolor{kiakzent}{rgb}{0.0,0.314,0.608}   % LUH-Blau
% Prompt-Box:   \begin{kiprompt}[Titel] ... \end{kiprompt}
\newtcolorbox{kiprompt}[1][Prompt]{%
    enhanced, breakable, sharp corners,
    colback=kiakzent!6, colframe=kiakzent, boxrule=0pt, leftrule=2.5pt,
    left=6pt, right=6pt, top=4pt, bottom=4pt,
    fonttitle=\bfseries\footnotesize, coltitle=kiakzent, colbacktitle=kiakzent!6,
    title={#1}, attach title to upper={\par\smallskip}, fontupper=\small,
    before skip=6pt, after skip=4pt,
}
% Antwort-Box:  \begin{kiantwort}[Titel] ... \end{kiantwort}
\newtcolorbox{kiantwort}[1][Antwort]{%
    enhanced, breakable, sharp corners,
    colback=black!3, colframe=black!45, boxrule=0pt, leftrule=2.5pt,
    left=6pt, right=6pt, top=4pt, bottom=4pt,
    fonttitle=\bfseries\footnotesize, coltitle=black!60,
    title={#1}, attach title to upper={\par\smallskip}, fontupper=\small,
    before skip=4pt, after skip=10pt,
}

%=== Grafiken ==========================================
\usepackage{graphicx}
\graphicspath{{figures/}}   % Bilderpfad
\usepackage{float}
\usepackage{wrapfig}
\usepackage{caption}
\usepackage{subcaption}
\usepackage[section]{placeins}

%=== Mathematik ========================================
\usepackage{amsmath}
\usepackage{amsfonts}
\usepackage{amssymb}

%=== Listen ============================================
\usepackage{enumitem}       % Flexible Listenabstände (ersetzt paralist)
\setlist{nosep}             % Kompakte Listen global (kein zusätzlicher vertikaler Abstand)

%=== Einheiten und Symbole =============================
\usepackage{siunitx}
\sisetup{locale=DE}         % Deutsche Formatierung
\usepackage[right]{eurosym}

%=== Typographie =======================================
\clubpenalty = 10000
\widowpenalty = 10000
\displaywidowpenalty = 10000
\tolerance = 10000
\hbadness = 10000
\emergencystretch = 3em
\hfuzz = 1pt

%=== Inhaltsverzeichnis ================================
\RedeclareSectionCommand[tocpagenumberwidth=2.5em]{chapter}

%=== Literaturverzeichnis (biblatex) ===================
\usepackage[
    backend=biber,
    style=numeric,
    sorting=none,
    doi=true,
    url=true,
    isbn=true,
    maxbibnames=99,
    giveninits=true,
]{biblatex}

\addbibresource{99_bibliography.bib}

% DOI als klickbaren Link formatieren
\DeclareFieldFormat{doi}{%
    \mkbibacro{DOI}\addcolon\space
    \ifhyperref
        {\href{https://doi.org/#1}{\nolinkurl{#1}}}
        {\nolinkurl{#1}}%
}

% Warnungen bei fehlenden DOI/URL/ISBN
\AtEveryBibitem{%
    % Artikel: DOI oder URL erforderlich
    \ifentrytype{article}{%
        \iffieldundef{doi}{%
            \iffieldundef{url}{%
                \PackageWarning{biblatex}{Fehlende DOI/URL für Artikel: \thefield{entrykey}}%
            }{}%
        }{}%
    }{}%
    % Bücher: ISBN oder DOI erforderlich
    \ifentrytype{book}{%
        \iffieldundef{isbn}{%
            \iffieldundef{doi}{%
                \PackageWarning{biblatex}{Fehlende ISBN/DOI für Buch: \thefield{entrykey}}%
            }{}%
        }{}%
    }{}%
}
}%
}

% ============================================================================
% ERSTE TITELSEITE (immer angezeigt)
% ============================================================================

\thispagestyle{empty}

\begin{center}

\begin{tabular}{p{6cm}p{6cm}}
    \includegraphics[width=6cm]{figures/logo_iwes}  &
    \includegraphics[width=6cm]{figures/logo_luh}\\
    &  \vspace{1cm} \large\name  \\
    &  \vspace{0.5cm} \setstretch{1.0}\Large\textbf{\titel} \\
\end{tabular}

\vspace{1cm}

\vfill

\begin{tabular}{p{6cm}p{6cm}}
    &  \fbox{\parbox[t]{5cm}{\centering{\textbf{\projectNumber}}}}  \vspace{2cm} \\
    &  \universityName \\
    &   \instituteName \\
    &   \professorName \\
    &   \instituteStreet \\
    &   \instituteCity
\end{tabular}

\end{center}

\cleardoublepage

% ============================================================================
% ZWEITE TITELSEITE
% ============================================================================

\thispagestyle{empty}

\begin{center}

\begin{tabular}{p{6cm}p{6cm}}
    \includegraphics[width=6cm]{figures/logo_iwes}  &
    \includegraphics[width=6cm]{figures/logo_luh}\\
    &  \vspace{1cm} \large\name  \\
    &  \vspace{0.5cm} \setstretch{1.0}\Large\textbf{\titel} \\
\end{tabular}

\vspace{1cm}

% Typ-Beschreibung (automatisch gesetzt durch \setthesistype)
\begin{tabular}{p{6cm}p{6cm}}
    & {\thesisdesc} \vspace{0.5cm}\\
    \ifbool{showdegree}{%
        & {\large\textbf{\projectType}}%
    }{%
        & {\textbf{\thesismodule}} \vspace{1cm}%
    }
\end{tabular}

\vfill

% Informationen
\begin{tabular}{p{6cm}p{6cm}}
    \raggedleft Studiengang: & \studiengang \\
    \raggedleft Matrikelnummer: &  \matrikelnummer \\
    \raggedleft vorgelegt am: &  \today \vspace{1.5cm} \\

    \raggedleft Prüfer: &  \textbf{\pruefer} \\
    & \fpruefer \vspace{0.5cm} \\

    \ifbool{showcoexaminer}{%
        \raggedleft Zweitprüfer: &  \textbf{\zweitpruefer} \\
        & \fzweitpruefer \vspace{0.5cm} \\
    }{}

    \raggedleft Betreuer: & \textbf{\betreuer} \\
    & \fbetreuer
\end{tabular}

\end{center}

\clearpage


\pagenumbering{Roman}
\setcounter{page}{2}

%=== Vorspann ==========================================
% Sperrvermerk und Eidesstattliche Erklärung
% Sperrvermerk ist optional - den Abschnitt auskommentieren wenn nicht benötigt

%=== Sperrvermerk (optional) ===========================
% Hinweis: \companyname in 00_main.tex setzen (im Studierenden-Konfigurationsbereich)
% Diesen Abschnitt auskommentieren wenn kein Sperrvermerk benötigt wird

\thispagestyle{plain}
\vspace*{9mm}
\section*{Sperrvermerk}
\label{Sperrvermerk}
\vspace*{11mm}

Die vorliegende Ausarbeitung mit dem Titel
%
\vspace*{5mm}
\begin{center}
\Large\emph{\titel}
\end{center}
\vspace*{5mm}
%
beinhaltet vertrauliche Daten und Informationen von \companyname.\\

Die Ausarbeitung wird nur dem Erst- und Zweitgutachter sowie befugten Mitgliedern des Prüfungsausschusses der Leibniz Universität Hannover zu dem Zweck der Begutachtung und Bewertung in schriftlicher und digitaler Form zugänglich gemacht.\\

Die Weitergabe oder Veröffentlichung des Inhaltes dieser Arbeit im Gesamten oder in Teilen, sowie das Anfertigen weiterer Kopien oder Abschriften auch in digitaler Form, ist nicht gestattet.\\

Ausnahmen bedürfen der schriftlichen Genehmigung durch \companyname.

\clearpage

%=== Eidesstattliche Erklärung =========================

\lhead{\slshape{Eidesstattliche Erklärung}}
\vspace*{9mm}
\section*{Eidesstattliche Erklärung}
\label{Erklaerung}
\vspace*{11mm}

Hiermit erkläre ich, \name, dass ich die vorliegende Ausarbeitung mit dem Titel
%
\begin{center}
\Large{\emph{\titel}}
\end{center}
%
ohne unzulässige Hilfe Dritter und ohne Benutzung anderer als der angegebenen Hilfsmittel angefertigt habe. Die aus fremden Quellen direkt oder indirekt übernommenen Gedanken sind als solche kenntlich gemacht. Methoden der künstlichen Intelligenz habe ich nicht verwendet oder entsprechende Passagen kenntlich gemacht. Art und Umfang des Einsatzes von KI-Tools sind in Anhang~\ref{cha:KI-Dokumentation} dokumentiert.% <- Satz entfernen, wenn keine KI-Tools verwendet wurden

Ich habe keine kommerzielle Beratung in Anspruch genommen. Dritte haben von mir weder unmittelbar noch mittelbar geldwerte Leistungen für Tätigkeiten erhalten, die im Zusammenhang mit dem Inhalt der vorgelegten Arbeit stehen.

Die Arbeit wurde bisher weder im Inland noch im Ausland in gleicher oder ähnlicher Form eingereicht und ist als Ganzes auch noch nicht veröffentlicht.

\vspace*{15mm}
\name, Hannover den \today

\clearpage
     % Enthält Sperrvermerk (optional) + Eidesstattliche Erklärung
% Aufgabenstellung
\lhead{\slshape{Aufgabenstellung}}
\pdfbookmark[1]{Aufgabenstellung}{Aufgabe}
\vspace*{9mm}
\section*{Aufgabenstellung}
\vspace*{11mm}

% Hinweis: Hier die eigentliche Aufgabenstellung einfügen
% Alternativ: PDF der Aufgabenstellung mit \includepdf einbinden

Lorem ipsum dolor sit amet, consetetur sadipscing elitr, sed diam nonumy eirmod tempor invidunt ut labore et dolore magna aliquyam erat, sed diam voluptua. At vero eos et accusam et justo duo dolores et ea rebum. Stet clita kasd gubergren, no sea takimata sanctus est Lorem ipsum dolor sit amet.

Duis autem vel eum iriure dolor in hendrerit in vulputate velit esse molestie consequat, vel illum dolore eu feugiat nulla facilisis at vero eros et accumsan et iusto odio dignissim qui blandit praesent luptatum zzril delenit augue duis dolore te feugait nulla facilisi.

Ut wisi enim ad minim veniam, quis nostrud exerci tation ullamcorper suscipit lobortis nisl ut aliquip ex ea commodo consequat.

\clearpage

% Kurzfassung (Abstract)
\lhead{\slshape{Kurzfassung}}
\pdfbookmark[1]{Kurzfassung}{Abstrakt}
\vspace*{0mm}
\section*{Kurzfassung}
\phantomsection

% Deutsche Kurzfassung
Für das Modul \textit{Wissenschaftliches Arbeiten im Ingenieurwesen} wird eine wissenschaftliche Ausarbeitung zu einem vorgegebenen Thema gefordert. In aller Regel ist als Bestandteil der Aufgabenstellung eine Literaturrecherche zum vorgegebenen Thema durchzuführen. Aus dieser ist der Stand der Wissenschaft und Technik des vorgegebenen Themas abzuleiten und schriftlich zusammenzufassen. Abhängig von der Aufgabenstellung kann zusätzlich die programmiertechnische Umsetzung von Modellen, deren Anwendung und die Dokumentation von Berechnungsergebnissen gefordert werden. Dieses Dokument stellt einen Leitfaden des Instituts für Windenergiesysteme der Leibniz Universität Hannover dar, der beim Anfertigen der schriftlichen Ausarbeitung unbedingt zu beachten ist. \\

\vspace*{9mm}
\section*{Abstract}

% Englische Kurzfassung
For the module \textit{Scientific Work in Engineering}, a scientific report on a given topic is required. A literature research on the given topic usually needs to be carried out as part of the report. From this, the state of the art in science and technology of the given topic must be derived and summarized in writing. Depending on the task, the programming implementation of models, their application and the documentation of calculation results may also be required. This document is a guideline from the Institute for Wind Energy Systems at Leibniz University Hannover, which must be followed when preparing the written report. \\

\clearpage


%=== Inhaltsverzeichnis ================================
\pdfbookmark[1]{Inhaltsverzeichnis}{toc}
\lhead{\slshape{Inhaltsverzeichnis}}
\tableofcontents
\cleardoublepage

%=== Hauptteil =========================================
\pagestyle{scrheadings}
\pagenumbering{arabic}
\setcounter{page}{1}

\chapter{Aufbau und Gliederung der Arbeit}
\label{cha:Aufbau}

Der Aufbau, die Gliederung und der Umfang einer wissenschaftlichen Arbeit variiert ja nach Art und Thema der Arbeit. Das folgende Gliederungsschema hat sich allerdings bewährt und sollte im Wesentlichen eingehalten werden:
%
\begin{itemize}
	\item Deckblatt und Titelseite
%	\item	\emph{Sperrvermerk}
	\item Eidesstattliche Erklärung
	\item Aufgabenstellung
	\item Kurzfassung und Abstract
	\item Inhaltsverzeichnis
	\item Einleitung
	\item Grundlagen
	\item Hauptteil
	\item Zusammenfassung und Ausblick
	\item \textit{Abbildungs-, Tabellen-, Symbol- und Abkürzungsverzeichnis}
	\item Literaturverzeichnis
	\item \textit{Anhänge}
\end{itemize}
%
Die kursiv gesetzten Einträge sind dabei als optional zu betrachten. In den folgenden Abschnitten wird kurz auf die einzelnen Gliederungspunkte eingegangen.

Es ist grundsätzlich darauf zu achten, dass nicht mehr als drei Gliederungsebenen verwendet werden, also Kapitel, Unterkapitel und Unter-Unterkapitel. Wenn weitere Gliederungsebenen erforderlich werden, ist das grundsätzliche Konzept der Ausarbeitung zu überdenken und entsprechend abzuändern.

\section{Deckblatt und Titelseite}
\label{sec:Titelseite}

Das Deckblatt und die Titelseite sind die ersten beiden Seiten dieses Leitfadens bzw. dieser Vorlage. Sie müssen entsprechend gestaltet werden und alle darauf aufgeführten Informationen enthalten. Auf dem Deckblatt sind das das Logo des Instituts, das Logo der Leibniz Universität Hannover, der Name der Autorin bzw. des Autors, der Titel der Arbeit, die Inventar-Nummer der Arbeit (diese wird kurz vor Abgabe der Arbeit vom Institut vergeben und ist in die Box zu setzen) sowie die Anschrift des Instituts. Auf der Titelseite sind der Passus \textit{Ausarbeitung für das Modul Wissenschaftliches Arbeiten im Ingenieurwesen}, der Studiengang, die Matrikelnummer, das Abgabedatum sowie der Prüfer und die Betreuerin bzw. der Betreuer gemäß dieser Vorlage zu ergänzen.

Das Deckblatt ist außen auf den Umschlag der gebundenen, gedruckten Ausgabe aufzudrucken. Die Titelseite bildet die erste Papierseite, die in den Umschlag eingebunden wird. Anschauungsexemplare können auf Anfrage am Institut für Windenergiesysteme eingesehen werden.

%\section{Sperrvermerk}
%\label{sec:Sperr}

%Bei wissenschaftlichen Arbeiten, die in Zusammenarbeit mit einem Unternehmen erstellt werden, kann ein Sperrvermerk der Arbeit gefordert werden. Der Sperrvermerk schützt dabei die vertraulichen Daten des Unternehmens gegenüber der Öffentlichkeit (z.B. Wettbewerber des Unternehmens) und ist als erste Seite unmittelbar nach dem Deckblatt in die Arbeit einzubinden. Optional kann der Sperrvermerk direkt auf dem Deckblatt platziert werden.

%Der Inhalt des Sperrvermerkes ist mit dem Betreuer abzustimmen.

\section{Eidesstattliche Erklärung}
\label{sec:Eid}

Eine eidesstattliche Erklärung ist ein wichtiger Bestandteil einer wissenschaftlichen Arbeit und muss der Arbeit vorangestellt werden. In der eidesstattlichen Erklärung versichert die Verfasserin bzw. der Verfasser, dass die vorliegende Arbeit eigenständig und ohne die Hilfe Dritter angefertigt wurde. Des Weiteren wird erklärt, dass alle verwendeten Quellen als solche kenntlich gemacht und im Literaturverzeichnis aufgeführt wurden. Außerdem wird durch den Verfasser versichert, dass die Arbeit in dieser Form weder in Deutschland noch im Ausland bereits veröffentlicht oder für ein Studium vorgelegt oder eingebracht wurde. Insbesondere wird auch bestätigt, dass künstliche Intelligenz zur Erstellung von Inhalten entweder nicht eingesetzt oder entsprechende Passagen kenntlich gemacht wurden. Wie der Einsatz von KI-Tools nach der Handreichung der Leibniz Universität Hannover dokumentiert werden kann, zeigt Anhang~\ref{cha:KI-Dokumentation}.

Der Wortlaut der eidesstattlichen Erklärung muss dem auf Seite \pageref{Erklaerung} entsprechen. Die Erklärung ist mit Ort, Datum sowie der Unterschrift der Verfasserin bzw. des Verfassers abzuschließen.

\section{Aufgabenstellung}
\label{sec:Aufgabe}

Die Aufgabenstellung gibt Auskunft über das Thema der Arbeit, beinhaltet die zu bearbeitenden Teilaufgaben und dokumentiert das Aus- und das spätestmögliche Abgabedatum sowie die Namen des Prüfers und der Betreuerin bzw. des Betreuers. Die Aufgabenstellung wird vom Institut für Windenergiesysteme ausgegeben. Es soll im Original hinter der Eidesstattlichen Erklärung eingebunden werden. Die zwei Seiten für die Aufgabenstellung dieser Vorlage sind als Platzhalter zu verstehen und sind durch die tatsächlich ausgegebene Aufgabenstellung zu ersetzen. 

\section{Kurzfassung und Abstract}
\label{sec:Abstract_DE}

Auf maximal einer halben Seite ist eine Kurzfassung der Arbeit voranzustellen. Die Kurzfassung soll der Leserin bzw. dem Leser einen Überblick über den Inhalt des Dokuments geben. Daher sollte sie eine kurze Einleitung ins Thema, die methodischen Ansätze, die zur Bearbeitung des Themas herangezogen wurden sowie die wesentlichen Ergebnisse der Bearbeitung beinhalten. Eine englischsprachige Kurzfassung (Abstract) ist nach der deutschsprachigen Kurzfassung zu ergänzen. Beide Kurzfassungen sollen insgesamt den Umfang einer Seite nicht überschreiten.

\section{Inhaltsverzeichnis}
\label{sec:Inhalt}

Das Inhaltsverzeichnis umfasst alle im Dokument vorhandenen Kapitel, Unterkapitel und Unter-Unterkapitel. Auch weitere Verzeichnisse (Abbildungs-, Tabellen-, Symbol- und Abkürzungsverzeichnisse) und ggf. Anhänge müssen im Literaturverzeichnis enthalten sein. 

Die Überschriften im Text und im Inhaltsverzeichnis müssen exakt übereinstimmen. Überschriften einzelner Gliederungspunkte dürfen sich nicht mit dem Gesamtthema der Arbeit und nicht mit Überschriften von Unterpunkten decken. Es ist darauf zu achten, dass aus den Überschriften der zu erwartende Inhalt hervorgeht. Allgemeingültige Kapiteltitel wie z. B. \textit{Grundlagen} oder \textit{Hauptteil} sind zu vermeiden.

Neben den durchnummerierten Abschnittsbezeichnungen ist die Seitenzahl des Abschnittsbeginns aufzuführen.

\section{Einleitung}
\label{sec:Einleitung}

Die Einleitung bildet das erste nummerierte Kapitel der Ausarbeitung. Sie leitet in das Thema ein, motiviert die Arbeit und leitet die Problemstellung ab. Die Einleitung dient dazu, das Thema in den übergeordneten Kontext zu setzen und ggf. eine Abgrenzung des Themas vorzunehmen. Den Abschluss der Einleitung bildet eine Beschreibung des Aufbaus der Ausarbeitung.

\section{Technisch-physikalische Grundlagen}
\label{sec:Grundlagen}

Das zweite Kapitel einer wissenschaftlichen Arbeit widmet sich den technisch-physikalischen Grundlagen, die ein fachlich gebildeter Leser kennen muss, um alle folgenden Kapitel nachvollziehen zu können. Es ist darauf zu achten, diesem Kapitel einen inhaltsbezogenen Titel zu geben. Ein allgemeiner Titel wie \textit{Technisch-physikalische Grundlagen} ist nicht zulässig.

\section{Hauptteil}
\label{sec:Hauptteil}

Im Hauptteil einer wissenschaftlichen Arbeit wird die Gesamtheit der bearbeiteten Problemstellung dargestellt. Die methodischen Ansätze sowie der Lösungsweg sind ebenso darzustellen wie die Ergebnisse der Untersuchungen. 

Generell ist auf eine klare Darstellung des Gedankengangs, eine der Nachvollziehbarkeit der Ausarbeitung dienliche Gliederung und folgerichtige Anordnung des Inhalts sowie auf eine genaue, verständliche, kurze und prägnante Formulierung zu achten.

Die Gliederung des Inhalts und damit die Reihenfolge der Darstellung müssen sich nach sachlichen Gesichtspunkten und nicht nach der zeitlichen Abfolge der Bearbeitung richten. Der Hauptteil kann mehrere thematisch getrennte Kapitel umfassen. Es ist darauf zu achten, alle Kapitel mit inhaltsbezogenen Titeln zu versehen. Ein allgemeingültiger Titel wie \textit{Hauptteil} ist nicht zulässig.

Durch die Beschreibung der durchgeführten Untersuchungen müssen die erzielten Ergebnisse reproduzierbar sein. Werden Ergebnisse oder Lösungen von anderen Quellen zitiert, so sind diese nur soweit zu erläutern, wie es für das Verständnis notwendig ist. Eine Herleitung kann hingegen entfallen. Allerdings sollten Herleitungen zu tatsächlich selbst verwendeten Methoden in jedem Fall eingebunden werden.

\section{Zusammenfassung und Ausblick}
\label{sec:Zusammenfassung}

Eine wissenschaftliche Arbeit wird durch eine Zusammenfassung abgeschlossen. Die Zusammenfassung hat klare Aussagen über die verwendeten Verfahren, die erhaltenen Ergebnisse und ihre Bedeutung, die Grenzen der Gültigkeit und die Fortschritte gegenüber dem Kenntnisstand zu Beginn der Arbeit sowie ggf. die Anwendbarkeit der Ergebnisse zu enthalten. 

Des Weiteren ist in der Zusammenfassung in geeigneter Weise auf den Eigenanteil der Arbeit hinzuweisen und eventuell ein Ausblick auf mögliche vertiefende Forschungsarbeiten oder weiterführende Untersuchungen der Thematik zu geben.

Eine Leserin bzw. ein Leser, der auf der Suche nach Informationen zu einem bestimmten Thema ist, sollte nach Durchsicht der Kurzfassung, der Zusammenfassung und der Einleitung einen guten Überblick über die gesamte Arbeit bekommen haben. Daher werden in der Zusammenfassung in der Regel Aussagen wiederholt, die im Hauptteil der Arbeit bereits ausführlich beschrieben worden sind.

\section{Abbildungs, Tabellen-, Symbol- und Abkürzungsverzeichnis}
\label{sec:Abb-Tab-Symb-Abk-v}

Das Abbildungsverzeichnis hat den Zweck, alle Abbildungen bzw. deren Bildunterschriften mit Angabe der Seitenzahlen aufzulisten. Das Abbildungsverzeichnis ist optional, muss allerdings eingefügt werden, wenn mehr als drei Abbildungen im Dokument enthalten sind. Das Tabellenverzeichnis ist analog zum Abbildungsverzeichnis einzubinden, wenn mehr als drei Tabellen im Dokument vorhanden sind.

Das Tabellenverzeichnis hat den gleichen Zweck wie das Abbildungsverzeichnis und ist ebenfalls optional.

Das Symbolverzeichnis dient der Auflistung aller verwendeter Symbole. Dies ist vor allem für Arbeiten mit vielen Formeln hilfreich, da Formelzeichen als Symbole aufgefasst werden. Die Aufführung erfolgt tabellarisch und sollte Angaben zu den verwendeten Symbolen, ihrer Bedeutung und Einheit enthalten. Das Symbolverzeichnis ist alphabetisch zu sortieren. Griechische und römische Symbole sind in getrennten Blöcken aufzulisten.

Das Abkürzungsverzeichnis wird einer wissenschaftlichen Arbeit dann beigefügt, wenn im Fließtext viele nicht geläufige Abkürzungen verwendet werden. In tabellarischer Form werden sowohl die Abkürzungen als auch ihre Bedeutungen aufgelistet. Hierbei sind nur nicht im "`Duden"' verzeichnete Abkürzungen anzugeben. 

\section{Literaturverzeichnis}
\label{sec:Literatur}

Im Literaturverzeichnis werden alle zitierten Quellen aufgeführt. Dies hat in alphabetischer Reihenfolge des Nachnamens der Erstautorinnen bzw. Erstautoren zu erfolgen. Die Nummern der Quellen haben mit den im Text verwendeten Hinweisnummern auf die entsprechende Literatur oder Quelle übereinzustimmen.

Weitere Informationen bezüglich der Zitierweise werden in Abschnitt \ref{sec:Zitier} gegeben. Layoutbeispiele für verschiedene Quellenarten finden sich im Literaturverzeichnis am Ende des Dokuments. Der Stil des Literaturverzeichnisses sollte in jedem Fall übernommen werden.

\section{Anhänge}
\label{sec:Anhänge}

Anhänge in einer wissenschaftlichen Arbeit dienen der Darstellung von Inhalten wie Abbildungen, Tabellen, Herleitungen, Programmquelltexten, etc., die im Text keinen Platz finden oder dort unnötig sind und dennoch für die Nachvollziehbarkeit und die Reproduzierbarkeit der Arbeit von Bedeutung sind. Anhänge werden im Inhaltsverzeichnis aufgeführt oder durch ein separates Verzeichnis aufgelistet. Im Fließtext sollte auf die Bezug genommenen Anhänge verwiesen werden.


\clearpage
\chapter{Formales}
\label{cha:form_DE}

In den folgenden Abschnitten wird auf verschiedene formale Aspekte eingegangen, die im Wesentlichen befolgt werden sollten. Die Vorgaben zur äußeren Form, der Gliederungstiefe und der Zitierweise sind unbedingt einzuhalten.

\section{Sprache}
\label{sec:sprache}

Wissenschaftliche Arbeiten am IWES können in deutscher oder englischer Sprache verfasst werden, wobei es für englischsprachige Arbeiten eine gesonderte Vorlage bzw. einen gesonderten Leitfaden gibt. Die Wahl der Sprache ist mit der Betreuerin bzw. dem Betreuer der Arbeit abzustimmen. In jedem Fall ist die gültige Rechtschreibung und Grammatik anzuwenden.

Der Sprachstil sollte einer wissenschaftlichen Ausarbeitung angemessen sein. Es ist Berichtssprache zu verwenden. Umgangssprachliche Formulierungen sind zu vermeiden. Wissenschaftliches Schreiben zeichnet sich durch einen sehr effizienten Sprachstil aus. Dieser ist jederzeit sachlich und so knapp wie möglich zu halten. Unnötige Füllwörter und Übergangsformulierungen sollten vermieden werden. Jedes Wort, das zum Verständnis der Aussage nicht unbedingt erforderlich ist, sollte weggelassen werden. Dabei ist natürlich trotzdem darauf zu achten, dass ein ansprechender Lesefluss gegeben ist.

Um die Fehlerdichte zu minimieren, sollten Tools für eine effektive Rechtschreib- und Grammatikprüfung vor der Abgabe der Arbeit Verwendung finden. Auch sprachliche Mängel können zu einer Abwertung der Ausarbeitung führen.

\section{Äußere Form}
\label{sec:form_DE}

Wissenschaftliche Ausarbeitungen sind einseitig, einspaltig und im Blocksatz bei 1,5-fachem Zeilenabstand zu erstellen. Als Schriftart ist eine serifenlose Schrift zu verwenden. Die Schriftgröße muss zwischen 11 pt und 12 pt gewählt werden. Alternativ kann, wenn die Arbeit mittels \LaTeX{} verfasst wird, die Standartschrift von \LaTeX{} bei 11 pt verwendet werden. Diese wird durch die Pakete \textbackslash{}usepackage[T1]\{fontenc\} und \textbackslash{}usepackage[utf8]\{inputenc\} auf den deutschen Zeichenvorrat ausgeweitet.

Abgesehen vom Deckblatt und der Titelseite, die eigene Formatierungen aufweisen, sind die folgenden Seitenränder einzuhalten:

\begin{itemize}
\item 2,5 cm (oben, links und rechts)
\item 2,0 cm (unten)
\end{itemize}

Kapitel müssen mit einer neuen Seite beginnen. Bei Unterkapiteln und Unter-Unterkapiteln wird der Seitenumbruch dagegen variabel gehalten.

\section{Gliederungsebenen, Kopfzeilen und Seitennummerierungen}
\label{sec:nummerierung}

Es dürfen nicht mehr als drei nummerierte Gliederungsebenen verwendet werden, also Kapitel, Unterkapitel und Unter-Unterkapitel. Sollte eine vierte Gliederungsebene erforderlich werden, ist die Ausarbeitung zu restrukturieren, um die Limitierung auf drei Gliederungsebenen einzuhalten.

Ab dem Titelblatt ist eine Kopfzeile vorzusehen. Am linken Rand der Kopfzeile werden die Kapitelnummer und die Kapitelbezeichnung aufgeführt. Die Seitennummerierung ist im rechten Teil der Kopfzeile anzuordnen. 

Die Nummerierung der einführenden Seiten (Eidesstattliche Erklärung, Aufgabenstellung, Kurzfassung und Inhaltsverzeichnis) erfolgt in römischen Zahlen (I, II, III, ...). Die römische Zählung beginnt mit der Titelseite (erste Papierseite). Ab den Verzeichnissen (Abbildungsverzeichnis etc.) ist die römische Nummerierung fortlaufend fortzuführen. 

\begin{addmargin}[1cm]{1cm}
	\textit{\textbf{Beispiel:} Endet das Inhaltsverzeichnis mit der römischen Seitenzahl \textit{VI}, ist die Seitennummerierung bei den Verzeichnissen mit der römischen Seitenzahl \textit{VII} weiterzuführen.}
\end{addmargin}

Der eigentliche Hauptteil des Textes, also alle Kapitel von der Einleitung bis zur Zusammenfassung und dem Ausblick, wird, mit 1 beginnend, mit arabischen Zahlen nummeriert (1, 2, 3, ...). Die Zählung beginnt mit der ersten Seite der Einleitung.

\section{Formeln, Abbildungen und Tabellen}

Die Anforderungen für die Nutzung und Einbindung von Formeln, Abbildungen und Tabellen werden in gesonderten Kapiteln dargelegt, s. die Abschnitte \ref{cha:Formeln}, \ref{cha:Images_DE} und \ref{cha:Tabellen}. Dort werden auch Beispiele für unterschiedliche Arten der Einbindung dargestellt. Gemeinsame Anforderungen werden im Folgenden definiert.

Formeln, Abbildungen und Tabellen sind zu nummerieren. Hierfür ist eine kapitelinterne, fortlaufende Nummerierung zu wählen. Das bedeutet, dass die jeweiligen Nummern mit der Kapitelnummer beginnen und nach einem folgenden Punkt die Nummerierung fortlaufend erfolgt. Die nach dem Punkt befindliche fortlaufende Nummerierung ist für Formeln, Abbildungen und Tabellen jeweils unabhängig voneinander und muss in jedem Kapitel wieder bei 1 beginnen. Das Prinzip ist in Tabelle \ref{tab:Nummerierung} zusammengefasst. Die Gliederungsebene innerhalb eines Kapitels wird bei der Nummerierung nicht berücksichtigt.

\begin{table}[!htbp]
	\small
	\centering
	\caption[Nummerierungsbeispiele für Formeln, Abbildungen und Tabellen]{Nummerierungsbeispiele für Formeln, Abbildungen und Tabellen}
	\begin{tabular}{P{2.5cm}P{1.5cm}P{1cm}P{2cm}P{1cm}P{1.5cm}P{1cm}}
		\toprule
		\textbf{Kapitel} & \textbf{Formel} & \textbf{Nr.} & \textbf{Abbildung} & \textbf{Nr.} & \textbf{Tabelle} & \textbf{Nr.} \\
		\midrule
		Kapitel 1 
		& 1. Formel & 1.1 & 1. Abbildung & 1.1 & 1. Tabelle & 1.1 \\
		& 2. Formel & 1.2 & 2. Abbildung & 1.2 & 2. Tabelle & 1.2 \\
		& ... & ... & ... & ... & ... & ... \\
		Kapitel 2 
		& 1. Formel & 2.1 & 1. Abbildung & 2.1 & 1. Tabelle & 2.1 \\
		& 2. Formel & 2.2 & 2. Abbildung & 2.2 & 2. Tabelle & 2.2 \\
		& ... & ... & ... & ... & ... & ... \\
		Kapitel 3
		& 1. Formel & 3.1 & 1. Abbildung & 3.1 & 1. Tabelle & 3.1 \\
		& 2. Formel & 3.2 & 2. Abbildung & 3.2 & 2. Tabelle & 3.2 \\
		... & ... & ... & ... & ... & ... & ... \\
		\bottomrule
	\end{tabular}%
	\label{tab:Nummerierung}%
\end{table}%



\section{Verweise und Fußnoten}
\label{sec:Verweise}

Innerhalb eines Kapitels sind oftmals Verweise auf andere Bestandteile der Arbeit notwendig. Verweise sollen mittels der jeweiligen Abschnittsnummer vorgenommen werden. Will man z.B. auf den Inhalt des Abschnittes "`Aufgabenstellung"' verweisen, so hat dies mit der Abschnittsnummer, in dem Fall \ref{sec:Aufgabe}, zu erfolgen. 

\begin{addmargin}[1cm]{1cm}
	\textit{\textbf{Beispiel:} Hinweise zum Einbinden der Aufgabenstellung finden sich in Abschnitt \ref{sec:Aufgabe}.}
\end{addmargin}

Derartige Verweise sind über die entsprechenden Befehle oder Funktionen der Textverarbeitungsprogramme zu erstellen. Ein händisches Eintragen der Nummern führt mit hoher Wahrscheinlichkeit zu Fehlern im Enddokument. 

Wird auf eine Seite verwiesen, die nicht Bestandteil eines Kapitels ist, kann alternativ auch durch die Angabe der Seitennummer auf dieses verwiesen werden. Dies ist aber im Allgemeinen zu vermeiden.

Wird auf Formeln, Abbildungen, Tabellen oder Anhänge Bezug genommen, so hat dies mittels der Angabe der Formel-, Abbildungs-, Tabellen- oder Anhangsnummer zu erfolgen (siehe hierzu auch Tabelle \ref{tab:Nummerierung}).

Fußnoten\footnote{Dies ist ein Beispiel einer Fußnote.} werden dann eingesetzt, wenn eine Aussage zwar für das Verständnis des Lesers wichtig, jedoch im Fließtext überflüssig oder fehlplatziert ist. Fußnoten werden über das gesamte Dokument hinweg durchgehend nummeriert. Die Fußzeile wird durch einen etwa 6,5 cm langen horizontalen Strich abgesetzt.

\section{Symbole und Abkürzungen}
\label{sec:Abk}

Symbole (auch Formelzeichen) und Abkürzungen sind innerhalb des Fließtextes dort einzuführen, wo sie zum ersten Mal verwendet werden. Für Symbole wird darauf in Abschnitt \ref{cha:Formeln} noch einmal eingegangen. Abkürzungen sind in Klammern hinter dem ausgeschriebenen Ausdruck einzuführen. 

\begin{addmargin}[1cm]{1cm}
	\textit{\textbf{Beispiel:} Dies ist der Leitfaden des Instituts für Windenergiesysteme (IWES) zum Anfertigen wissenschaftlicher Ausarbeitungen.}
\end{addmargin}



\section{Zitierrichtlinien}
\label{sec:Zitier}

Grundsätzlich sind alle aus anderen Veröffentlichungen wörtlich übernommenen Satzteile, Sätze oder Abschnitte sowie übernommene Abbildungen oder Tabellen zu zitieren (sogenannte direkte Zitate). Wörtlich übernommene Textbestandteile sind in Anführungsstrichen zu setzen. Das Zitat muss in dem Fall die Seiten- und Zeilenzahlen beinhalten. Sinngemäß übernommene Aussagen, Abbildungen, etc. sind ebenfalls durch die Angabe der Quelle als solche kenntlich zu machen (indirekte Zitate). Direkte Zitate sind im Ingenieurwesen äußerst unüblich und sollten vermieden werden. Indirekte Zitate sind grundsätzlich zu bevorzugen.

Eindeutige Quellenangaben sind beim wissenschaftlichen Arbeiten unumgänglich. Dadurch wird es der Leserin bzw. dem Leser ermöglicht, diese zu finden und somit die Herkunft der wissenschaftlichen Aussagen zu überprüfen, sofern es sich um nicht selbständig entwickeltes Gedankengut handelt.

Als Zitierstil ist in wissenschaftlichen Arbeiten am IWES eine numerische Zitierweise anzuwenden. Dafür ist die Nummer der jeweiligen Quelle im Literaturverzeichnis in eckigen Klammern hinter die zitierte Aussage zu setzen. 

\begin{addmargin}[1cm]{1cm}
	\textit{\textbf{Beispiel:} Eine detaillierte Beschreibung von Regelungskonzepten für große Windenergiesysteme findet sich in \cite{buch}.}
\end{addmargin}

Die Aufführung der Quellen im Literaturverzeichnis erfolgt in alphabetischer Reihenfolge. Es ist wichtig, dass die Formatierung von Einträgen im Literaturverzeichnis durchgängig konsistent ist. Dabei gibt es unterschiedliche Formatierungsvorschriften für unterschiedliche Quellenarten. Grundsätzlich sind überwiegend zitierfähige Quellen zu verwenden. Als zitierfähig gelten insbesondere Artikel in Fachzeitschriften \cite{artikel}, Fachbücher \cite{buch}, Beiträge in Konferenzbänden \cite{inp}, Teile eines Sammelbandes \cite{in}, technische Dokumentationen \cite{manual}, veröffentlichte Berichte einer Hochschule oder anderen Forschungsinstitution \cite{bericht} sowie Dissertationen \cite{thesis}. Nicht zitierfähige Quellen sind grundsätzlich zu vermeiden. Zwar kann insbesondere für Abbildungen nicht immer auf Internetquellen \cite{internet} verzichtet werden, dennoch sollten vorzugsweise Originalquellen (z. B. die Quellen von Wikipedia-Artikeln) verwendet werden. Die oben zitierten Referenzen stellen Beispiele dar. Die zugehörige Formatierung innerhalb des Literaturverzeichnisses sollte eingehalten werden.

\begin{addmargin}[1cm]{1cm}
	\textit{\textbf{Hinweis:} Vorlesungsunterlagen gelten als nicht zitierfähige Quellen, da sie als Allgemeinwissen behandelt werden. Im Zweifel sollte die in Vorlesungsunterlagen angegebene Originalliteratur zitiert werden. Auch Internetquellen sind als nicht zitierfähig einzustufen, da sie sich fortlaufend verändern und ggf. nicht überprüft werden können.}
\end{addmargin}


\section{Abgabe}
\label{sec:Abgabe}

Eine am IWES erstellte wissenschaftliche Arbeit ist in gedruckter und gebundener Form in einem weißen Umschlagkarton eingebunden abzugeben, auf dem das Deckblatt farbig aufgedruckt ist. 

\begin{addmargin}[1cm]{1cm}
	\textit{\textbf{Hinweis:} Gedruckte Fassungen in Ringbindung werden nicht angenommen, da eine Inventarisierung dann nicht möglich ist.}
\end{addmargin}


Die digitale Kopie der Arbeit ist als PDF-Datei per E-Mail an den Prüfer und die Betreuerin bzw. den Betreuer zu übersenden. Weitere elektronische Daten, einschließlich aller für die Reproduzierbarkeit relevanten Dateien (wie Quellen, Ergebnisse, Tabellen, Excel-Berechnungen, Simulationsergebnisse usw.), sollten über \textit{SeaFile} abgegeben werden. Der entsprechende Upload-Link wird auf Nachfrage von der Betreuerin bzw. dem Betreuer zur Verfügung gestellt.

Die Anzahl der abzugebenden gedruckten Fassungen ist der Aufgabenstellung zu entnehmen.




\clearpage
\chapter{Formeln}
\label{cha:Formeln}

Bei der Herleitung von Berechnungs- und/oder Simulationsmodellen ist auf eine mathematische Formulierung zurückzugreifen. Hierfür werden in der Regel Formeln benötigt. Formeln sind als Satzbestandteile zu behandeln und können entweder in den Fließtext eingebunden oder abgesetzt mit Nummerierung dargestellt werden. Es ist zwingend darauf zu achten, dass alle Formelzeichen (Symbole) im Fließtext eingeführt werden, und zwar an der Stelle, wo sie zum ersten Mal verwendet werden.

\begin{addmargin}[1cm]{1cm}
	\textit{\textbf{Beispiel:} Eine in den Fließtext eingebundene Formel kann z. B. der Satz des Pythagoras sein, der durch die Gleichung $a^2+b^2=c^2$ ausgedrückt wird, wobei $a$ und $b$ die Katheten und $c$ die Hypotenuse eines rechtwinkligen Dreiecks sind.}
\end{addmargin}

Der Satz des Pythagoras kann auch als abgesetzte und nummerierte Formel in der Form

\begin{equation}
	a^2+b^2=c^2
	\label{eq:pythagoras_DE}
\end{equation}

dargestellt werden. Auf eine in den Fließtext eingebettete Gleichung kann im Text nicht verwiesen werden. Daher ist darauf zu achten, dass Gleichungen, auf die Bezug genommen werden soll, in abgesetzter Form eingebunden werden. Bei einem Verweis auf eine Gleichung ist die Gleichungsnummer in runden Klammern zu referenzieren. 

\begin{addmargin}[1cm]{1cm}
	\textit{\textbf{Beispiel:} Gleichung (\ref{eq:pythagoras_DE}) zeigt den Satz des Pythagoras.}
\end{addmargin}

In einigen Fällen kann es zweckmäßig sein, einen Satz an Gleichungen mit ausgerichteten Gleichheitszeichen einzubinden, in der jede einzelne Gleichung eine eigene Nummer bekommt. Als Beispiel seien die allgemein bekannten binomischen Formeln in der Form

\begin{align}
(a+b)^2 & = a^2 + 2ab + b^2 \, , \label{eq:bin-1_DE} \\
(a-b)^2 & = a^2 - 2ab + b^2 \, , \label{eq:bin-2_DE} \\
(a+b)(a-b) & = a^2 - b^2 \, . \label{eq:bin-3_DE} 
\end{align}

Hierin sind $a$ und $b$ beliebige Variablen. Wie oben bereits erwähnt wurde, müssen Gleichungen immer als Satzbestandteil behandelt werden. Dies wird durch die Satzverbindungs- und Satzendezeichen (Kommas und Punkt) am Ende der Gleichungen (\ref{eq:bin-1_DE}), (\ref{eq:bin-2_DE}) und (\ref{eq:bin-3_DE}) deutlich. Matrizen sind auch häufig verwendete Optionen. Die Einheitsmatrix $\mathbf{I}$ wäre in abgesetzter Form beispielsweise durch den Ausdruck

\begin{equation}
	\mathbf{I} = 
	\left[ \begin{matrix}
		1 & 0 & 0 \\
		0 & 1 & 0 \\
		0 & 0 & 1 \\
	\end{matrix} \right]
\end{equation}

einzuführen. Grundsätzlich sind Vektoren durch fettgedruckte Kleinbuchstaben, Tensoren und Matrizen durch fettgedruckte Großbuchstaben zu kennzeichnen. Ausnahmen können u. U. griechische Buchstaben bilden, da mechanische Spannungen und Verzerrungen, obwohl sie tensorielle Größen sind, häufig durch die kleinen griechischen Buchstaben $\boldsymbol{\sigma}$ und $\boldsymbol{\varepsilon}$ ausgedrückt werden.

Wenn lange Gleichungen oder Gleichungen mit Umformungen so lang werden, dass sie nicht in eine Zeile passen, sollten sie mit ausgerichtetem Gleichheitszeichen auf mehrere Zeilen aufgeteilt werden. Schreibt man die erste binomische Formel, s. Gleichung (\ref{eq:bin-1_DE}),  mit Zwischenschritten, wird daraus beispielsweise der Ausdruck

\begin{equation}
	\begin{split}
		(a+b)^2 &= (a+b)(a+b)\\
		&= a^2 + ab + ab + b^2 = a^2 + 2ab + b^2 \,.
	\end{split}
\end{equation}

Die Einbindung von Brüchen, Produkten oder Integralen seien ebenfalls beispielhaft dargestellt durch die Beziehungen

\begin{equation}
	\binom{n}{k} = \frac{\displaystyle\prod_{i=1}^n i}{\displaystyle\prod_{i=1}^k i\cdot
		\prod_{i=1}^{n-k} i}
\end{equation}

und

\begin{equation}
	\int_0^5\limits \left(x^2 +x+1\right) \mathrm{d}x = \left. \frac{1}{3} x^3 + \frac{1}{2} x^2 + x + c \; \right|_0^5 = \frac{355}{6} + c \, ,
\end{equation}

wobei $x$ eine Variable und $c$ eine Integrationskonstante ist.
\clearpage
\chapter{Abbildungen}
\label{cha:Images_DE}

Die grafische Darstellung von Ergebnissen oder Zusammenhängen ist in wissenschaftlichen Arbeiten ein wichtiges Mittel. Abbildungen können auf verschiedene Art und Weise in den Fließtext einer Arbeit integriert werden. Am häufigsten wird die zentrierte und vom Text abgesetzte Einbindung verwendet. Exemplarisch ist Abbildung \ref{fig:centered_DE} auf diese Weise eingebunden.

\begin{figure}[!hb]
	\centering
	\includegraphics[width=0.4\linewidth, height=0.4\textheight, keepaspectratio]{figures/example-image.pdf}
	\caption[Zentrierte und vom Text abgesetzte Abbildung.]{Zentrierte und vom Text abgesetzte Abbildung.}
	\label{fig:centered_DE}
\end{figure}

Abbildungen müssen immer in einer Größe eingebunden werden, dass Texte und Details lesbar bzw. erkennbar bleiben. Idealerweise orientieren sich Schriftgrößen in Abbildungen an der Schriftgröße des Fließtextes. Ggf. kann sie eine Schriftgröße kleiner ausfallen als der Fließtext. 

\begin{wrapfigure}[12]{l}{0.41\textwidth}
	\begin{minipage}{0.95\linewidth}
		\vspace{-\baselineskip}
		\includegraphics[width=\textwidth, keepaspectratio]{figures/example-image.pdf}
		\caption[Umflossene Abbildung.]{Umflossene Abbildung.}
		\label{fig:wrapped-left_DE}
	\end{minipage}\hfill
\end{wrapfigure}

Werden ursprünglich farbige Abbildungen schwarzweiß gedruckt, muss eine ausreichende Unterscheidbarkeit der Zusammenhänge gewährleistet bleiben. Dies kann z. B. durch Verwendung unterschiedlicher Linienarten (durchgehende, gestrichelte oder gepunktete Linien) in Diagrammen mit mehr als einer Kurve umgesetzt werden. 

Mit Text umflossene Abbildungen, s. Abbildung \ref{fig:wrapped-left_DE}, sind grundsätzlich möglich. Diese eignen sich allerdings nur für kleine oder vertikal ausgerichtete Abbildungen, die horizontal nicht mehr als die Hälfte der Textbreite in Anspruch nehmen. Bei umflossenen Abbildungen ist die Anpassung der Abstände zum Text häufig schwierig, insbesondere in vertikaler Richtung. Im Abbildung \ref{fig:wrapped-left_DE} ist beispielsweise zu erkennen, dass die Beschriftung nicht genau am Seitenende endet. Daher sollte im Allgemeinen auf umflossene Abbildungen verzichtet werden.

Es gibt die Möglichkeit, mehrere schmale Abbildungen innerhalb einer Abbildung darzustellen, s. Abbildung \ref{fig:subfigure_DE}. Diese Anordnung nennt man auch \textit{Subfigure}-Konfiguration. In diesem Fall ist jede Teilabbildung (Subfigure) mit römischen Buchstaben, die in runden Klammern mittig unterhalb jeder Teilabbildung zu setzen sind, zu nummerieren. Die Abbildungsbeschriftung muss den Inhalt beider Teilabbildungen enthalten und auf die Buchstabennummerierung verweisen, s. auch die Abbildungsbeschriftung zu Abbildung \ref{fig:subfigure_DE}. Eine Subfigure-Konfiguration ist zu zentrieren, wobei üblicherweise versucht wird, die vollständige Textbreite zu nutzen.

\begin{figure}[!htbp]
	\centering
	\begin{subfigure}[b]{0.4\linewidth}
		\centering
		\includegraphics[width=\linewidth, keepaspectratio]{figures/example-image.pdf}
		\caption{}
	\end{subfigure}
	\hspace{0.5cm}
	\begin{subfigure}[b]{0.4\linewidth}
		\centering
		\includegraphics[width=\linewidth, keepaspectratio]{figures/example-image.pdf}
		\caption{}
	\end{subfigure}
	\caption[Zentrierte und vom Text abgesetzte Abbildung mit zwei Unterabbildungen, die sogenannte Subfigure-Konfiguration. (a) Beispielbild links angeordnet. (b) Beispielbild rechts angeordnet.]{Zentrierte und vom Text abgesetzte Abbildung mit zwei Unterabbildungen, die sogenannte Subfigure-Konfiguration. (a) Beispielbild links angeordnet. (b) Beispielbild rechts angeordnet. Ein bisschen mehr Text, der im Abbildungsverzeichnis entfällt.}
	\label{fig:subfigure_DE}
\end{figure}

Eine Abbildungsunterschrift muss die Abbildung so beschreiben, dass auch ohne den Fließtext die Abbildung für sich stehend verständlich ist. Wenn z. B. bei Graphen die Achsen mit Symbolen (Formelzeichen) beschriftet werden, sind die Symbole auch in der Abbildungsbeschriftung einzuführen, s. auch Abschnitt \ref{cha:Formeln}. Bei langen Abbildungsbeschriftungen bietet \LaTeX{} die zweckmäßige Möglichkeit, eine Kurzform in das Abbildungsverzeichnis zu übernehmen, s. auch den Eintrag im Abbildungsverzeichnis für Abbildung \ref{fig:subfigure_DE}.

Die Abbildungsbeschriftung ist unterhalb der Abbildung zu platzieren und im Blocksatz mit hängendem Einzug zu formatieren. Der hängende Einzug ist so einzustellen, dass der Text linksbündig nach der Abbildungsnummerierung ausgerichtet ist, s. auch die Abbildungsbeschriftung zu Abbildung \ref{fig:subfigure_DE}.

Ist eine Abbildung nicht selbst erstellt, ist die Quelle in der Abbildungsbeschriftung zu referenzieren. Anderenfalls handelt es sich um ein Plagiat. Eine Abbildung ohne Quellenangabe wird als selbst erstellt aufgefasst. Es ist nicht gesondert in der Abbildungsbeschriftung zu vermerken, z. B. durch den Zusatz " `selbst erstellt" '. 

Grundsätzlich ist auf jede Abbildung im Text zu verweisen. Abbildungen ohne Verweis sind zu entfernen. Der Inhalt bzw. der Einbindungs-Zweck jeder Abbildung ist im Fließtext zumindest kurz zu beschreiben bzw. zu benennen. 
\clearpage
\chapter{Tabellen}
\label{cha:Tabellen}

Tabellen sind neben Abbildungen ein zweckmäßiges Mittel, um Inhalte übersichtlich darzustellen oder Daten zu integrieren. Sie sind zentriert und vom Text abgesetzt einzubinden. Von umflossenen Tabellen ist abzusehen.

Wie auch Abbildungen sind Tabellen mit Beschriftungen zu versehen. Die Tabellenbeschriftung ist oberhalb der Tabelle anzuordnen, s. z. B. Tabelle \ref{tab:Beispieltab-1_DE}. Sie wird zentriert und im Blocksatz mit hängendem Einzug gesetzt, so dass die eigentliche Beschriftung linksbündig nach der Nummerierung beginnt, s. auch die Regelungen zu Abbildungen in Abschnitt \ref{cha:Images_DE}. 

\begin{table}[!htbp]
	\centering
	\caption[Zweispaltige Beispieltabelle mit unterschiedlichen Spaltenausrichtungen.]{Zweispaltige Beispieltabelle mit unterschiedlichen Spaltenausrichtungen (linke Spalte rechtsbündig, rechte Spalte linksbündig). Die Beschriftung hat einen hängenden Einzug.}
	\begin{tabular}{rl}
		\toprule
		\textbf{Arbeitsschritt} & \textbf{Teilarbeitsschritt} \\
		\midrule
		Einfacher Arbeitsschritt & Irgendetwas Einfaches \\
		& Noch irgendetwas Einfaches \\
		&\\
		Etwas komplizierterer Arbeitsschritt & Vorbereitende Arbeiten \\
		& Irgendetwas Kompliziertes \\
		& Noch irgendetwas Kompliziertes \\
		& Abschließende Arbeiten, die noch komplizierter sind \\
		\bottomrule
	\end{tabular}%
	\label{tab:Beispieltab-1_DE}%
\end{table}%

Die Beschriftung muss den Tabelleninhalt eindeutig beschreiben, so dass die Tabelle auch ohne Fließtext für sich stehend verständlich ist. Eventuell verwendete Symbole (Formelzeichen) sind in der Bildbeschriftung einzuführen. Bei langen Tabellenbeschriftungen gibt \LaTeX{} wie auch bei Abbildungen die Möglichkeit, eine Kurzform für das Tabellenverzeichnis zu definieren.

\begin{addmargin}[1cm]{1cm}
	\textit{\textbf{Hinweis:} Auf jede Tabelle muss im Fließtext Bezug genommen werden. Nicht referenzierte Tabellen sind zu entfernen.}
\end{addmargin}

Wird eine Tabelle aus einer anderen Quelle übernommen, ist diese in der Tabellenbeschriftung zu referenzieren. Es ist allerdings zu vermeiden, Tabellen aus anderen Quellen unverändert zu übernehmen, da sonst ein einheitlicher Tabellenstil nicht möglich ist. Besser ist es, nur den Inhalt einer Tabelle aus einer Quelle zu übernehmen und diesen in eine selbst gestaltete Tabelle zu integrieren. Auch in diesem Fall ist die Quelle zu referenzieren, z. B. in der Form " `Inhalte aus [XX]" ' (wobei \textit{XX} ein Platzhalter für die Referenznummer ist).

Tabellen dürfen nicht über die Textbreite hinausragen. In breiteren Tabellen kann in Grenzen die Schriftgröße verkleinert werden, s. Tabelle \ref{tab:Beispieltab-2_DE}. Es ist aber in jedem Fall darauf zu achten, dass die Lesbarkeit gewährleistet ist. Die Schriftgröße in Tabelle \ref{tab:Beispieltab-2_DE} ist grenzwertig. Kleiner sollte die Schriftgröße nicht werden. Wenn die Tabellen sehr breit sind, sollten sie im Querformat eingebunden werden.

\begin{table}[t]%[!htbp]
	\tiny
	\centering
	\caption[Siebenspaltige Beispieltabelle mit kleiner Schrift.]{Siebenspaltige Beispieltabelle mit kleiner Schrift, damit die Tabelle in die Textbreite passt. Kleiner als in diesem Beispiel sollte die Schriftgröße innerhalb einer Tabelle nicht sein.}
	\begin{tabular}{ccccccc}
		\toprule
		\textbf{Überschrift 1} & \textbf{Überschrift 2} & \textbf{Überschrift 3} & \textbf{Überschrift 4} & \textbf{Überschrift 5} & \textbf{Überschrift 6} & \textbf{Überschrift 7} \\
		\midrule
		Inhalt zu Überschrift 1 & Inhalt zu Überschrift 2 & Inhalt zu Überschrift 3 & Inhalt zu Überschrift 4 & ... & ... & ... \\
		Inhalt zu Überschrift 1 & Inhalt zu Überschrift 2 & Inhalt zu Überschrift 3 & Inhalt zu Überschrift 4 & ... & ... & ... \\
		Inhalt zu Überschrift 1 & Inhalt zu Überschrift 2 & Inhalt zu Überschrift 3 & ... & ... & ... & ... \\
		Inhalt zu Überschrift 1 & Inhalt zu Überschrift 2 & ... & ... & ... & ... & ... \\
		Inhalt zu Überschrift 1 & ... & ... & ... & ... & ... & ... \\
		... & ... & ... & ... & ... & ... & ... \\
		\bottomrule
	\end{tabular}%
	\label{tab:Beispieltab-2_DE}%
\end{table}%

Es gibt unzählige Möglichkeiten, wie eine Tabelle im Detail gestaltet werden kann. Für ein einheitliches Layout von Tabellen sollte auf vertikale Linien verzichtet werden, wenn sie nicht zwingend erforderlich sind. Oberhalb und unterhalb der Spaltenbeschriftungen und unterhalb der letzten Tabellenzeile ist eine horizontale Linie anzuordnen, wobei die oberste und unterste Linie etwas dicker sein soll, s. Tabelle \ref{tab:Beispieltab-1_DE} und Tabelle \ref{tab:Beispieltab-2_DE}. Zur vertikalen Abgrenzung können Leerzeilen eingefügt werden, s. Tabelle \ref{tab:Beispieltab-1_DE}. Alternativ können zusätzliche dünne horizontale Linien zur Unterteilung von Tabellenblöcken sinnvoll sein, s. Tabelle \ref{tab:Beispieltab-3_DE} und Tabelle \ref{tab:Beispieltab-4_DE}.

Die horizonalte Ausrichtung von Spalten (linksbündig, rechtsbündig oder zentriert) ist im Einzelfall abzuwägen. Die Tabellen \ref{tab:Beispieltab-1_DE} bis \ref{tab:Beispieltab-4_DE} zeigen die verschiedenen Möglichkeiten, auch in Kombination. Mitunter kann es zweckmäßig sein, Zeilen teilweise vertikal zu verbinden. Ein Beispiel dafür ist in Tabelle \ref{tab:Beispieltab-4_DE} dargestellt. Dort sind in den ersten beiden Spalten die Zellen vertikal verbunden, so dass in der dritten Spalte eine Zeilenaufteilung erkennbar ist. Die dünnen horizontalen Linien erhöhen in diesem Fall die Lesbarkeit der Tabelle.

\begin{table}[!hbpt]
	\tiny
	\centering
	\caption[Sechsspaltige Beispieltabelle mit kleiner Schrift.]{Sechsspaltige Beispieltabelle. Die dritte Spalte ist linksbündig, alle anderen Spalten sind zentriert.}
	\begin{tabular}{cclccc}
		\toprule
		\textbf{Zu\-ordnung} & \textbf{Pos.} & \textbf{Positionsbeschreibung} & \multicolumn{1}{c}{\textbf{GP [\EUR{}]}} & \multicolumn{1}{c}{\textbf{Prozent von GP}} & \multicolumn{1}{c}{\textbf{GP [\EUR{}]}} \\
		\midrule
		\multicolumn{1}{c}{\multirow{14}[0]{*}{A - Pos.}} & 3.5.10 & \multicolumn{1}{l}{Stb-Wand, C20/25; h = bis 19,05 m; d = 0,2 m herstellen} & \EUR{547.720,07} & \multicolumn{1}{c}{15,36\%} & \EUR{547.720,07} \\
		\multicolumn{1}{c}{} & 1.2.10 & \multicolumn{1}{l}{Vorhalten der Baustelle} & \EUR{420.308,58} & \multicolumn{1}{c}{27,15\%} & \EUR{968.028,65} \\
		\multicolumn{1}{c}{} & 4.2.10 & \multicolumn{1}{l}{Filigrandecke d = 0,05 m liefern und verlegen} & \EUR{312.057,83} & \multicolumn{1}{c}{35,90\%} & \EUR{1.280.086,48} \\
		\multicolumn{1}{c}{} & 4.8.20 & \multicolumn{1}{l}{Stb-Stützen, C30/37, b/d = 0,6/0,9, h = 18,75 m herstellen} & \EUR{226.309,35} & \multicolumn{1}{c}{42,25\%} & \EUR{1.506.395,83} \\
		\multicolumn{1}{c}{} & 4.2.20 & \multicolumn{1}{l}{Filigrandecke Maschinenhalle d = 0,05 m liefern u. verlegen} & \EUR{224.606,69} & \multicolumn{1}{c}{48,55\%} & \EUR{1.731.002,51} \\
		\multicolumn{1}{c}{} & 3.3.10 & \multicolumn{1}{l}{Bodenplatte, Stb C30/37, d = 25 cm herstellen} & \EUR{173.623,54} & \multicolumn{1}{c}{53,41\%} & \EUR{1.904.626,06} \\
		\multicolumn{1}{c}{} & 3.4.30 & \multicolumn{1}{l}{Stb-Stützen, C30/37, b/d = 0,6/2,8, h = 6,25 m herstellen} & \EUR{164.099,59} & \multicolumn{1}{c}{58,02\%} & \EUR{2.068.725,65} \\
		\multicolumn{1}{c}{} & 4.2.60 & \multicolumn{1}{l}{Aufbeton der Decke EG d = 0,35 m herstellen} & \EUR{151.737,34} & \multicolumn{1}{c}{62,27\%} & \EUR{2.220.462,99} \\
		\multicolumn{1}{c}{} & 2.1.20 & \multicolumn{1}{l}{Baugrubenaushub Bkl. 3-5, inkl. Abfuhr} & \EUR{134.779,63} & \multicolumn{1}{c}{66,05\%} & \EUR{2.355.242,63} \\
		\multicolumn{1}{c}{} & 3.1.20 & \multicolumn{1}{l}{Einzelfundamente Stb C20/25, b/l = 2,10/2,80; h = 1,50 m} & \EUR{126.460,84} & \multicolumn{1}{c}{69,60\%} & \EUR{2.481.703,47} \\
		\multicolumn{1}{c}{} & 4.6.10 & \multicolumn{1}{l}{Entwässerungsrinne: Breite 60 cm; Hmax = 1,50 m} & \EUR{115.367,81} & \multicolumn{1}{c}{72,83\%} & \EUR{2.597.071,28} \\
		\multicolumn{1}{c}{} & 3.2.10 & \multicolumn{1}{l}{Streifenfundamente Stb C20/25 b = 1,40 m; l = 10 m, 70 m; h = 1,50 m} & \EUR{91.621,95} & \multicolumn{1}{c}{75,40\%} & \EUR{2.688.693,23} \\
		\multicolumn{1}{c}{} & 1.1.10 & \multicolumn{1}{l}{Einrichten der Baustelle} & \EUR{88.494,12} & \multicolumn{1}{c}{77,88\%} & \EUR{2.777.187,35} \\
		\multicolumn{1}{c}{} & 4.2.50 & \multicolumn{1}{l}{Aufbeton Maschinenhalle Stb C20/25 d = 0,15 m herstellen} & \EUR{86.053,13} & \multicolumn{1}{c}{80,30\%} & \EUR{2.863.240,48} \\
		\midrule
		\multicolumn{1}{c}{\multirow{9}[0]{*}{B - Pos.}} & 4.7.20 & \multicolumn{1}{l}{Stb-Stützen, C30/37, b/d = 0,6/0,6, h = 18,75 m herstellen} & \EUR{83.945,37} & \multicolumn{1}{c}{82,65\%} & \EUR{2.947.185,85} \\
		\multicolumn{1}{c}{} & 4.4.10 & \multicolumn{1}{l}{Deckenunterzüge, Stb C35/45 , b/d = 0,6/0,8 m herstellen} & \EUR{72.332,54} & \multicolumn{1}{c}{84,68\%} & \EUR{3.019.518,39} \\
		\multicolumn{1}{c}{} & 1.3.10 & \multicolumn{1}{l}{Räumen der Baustelle} & \EUR{71.922,97} & \multicolumn{1}{c}{86,70\%} & \EUR{3.091.441,36} \\
		\multicolumn{1}{c}{} & 4.3.10 & \multicolumn{1}{l}{Spannbetonbinder C35/45 Hallendecke herstellen} & \EUR{65.125,35} & \multicolumn{1}{c}{88,52\%} & \EUR{3.156.566,71} \\
		\multicolumn{1}{c}{} & 3.1.10 & \multicolumn{1}{l}{Einzelfundamente Stb C20/25, b/l = 1,60 /2,80; h = 1,50 m} & \EUR{52.688,79} & \multicolumn{1}{c}{90,00\%} & \EUR{3.209.255,50} \\
		\multicolumn{1}{c}{} & 3.2.20 & \multicolumn{1}{l}{Streifenfundamente Stb C20/25, b/l = 0,6/4,8; h = 1,50 m} & \EUR{46.704,40} & \multicolumn{1}{c}{91,31\%} & \EUR{3.255.959,90} \\
		\multicolumn{1}{c}{} & 4.7.10 & \multicolumn{1}{l}{Stb-Stützen, C30/37, b/d = 0,6/0,6, h = 10,95 m herstellen} & \EUR{44.146,46} & \multicolumn{1}{c}{92,55\%} & \EUR{3.300.106,36} \\
		\multicolumn{1}{c}{} & 2.2.20 & \multicolumn{1}{l}{Trag.- und Frostschutzschicht liefern und herstellen} & \EUR{43.487,87} & \multicolumn{1}{c}{93,77\%} & \EUR{3.343.594,22} \\
		\multicolumn{1}{c}{} & 4.8.10 & \multicolumn{1}{l}{Stb-Stützen, C30/37, b/d = 0,6/0,9, h = 6,25 m herstellen} & \EUR{42.400,87} & \multicolumn{1}{c}{94,96\%} & \EUR{3.385.995,09} \\
		\bottomrule
	\end{tabular}%
	\label{tab:Beispieltab-3_DE}%
\end{table}%

\begin{table}[htbp!]
	\centering
	\caption[Dreispaltige Beispieltabelle mit Gruppenstruktur.]{Dreispaltige Beispieltabelle mit Gruppenstruktur.}
	\begin{tabularx}{\textwidth}{l l X}
		\toprule
		\textbf{Fachbereich} & \textbf{Kernaufgabe} & \textbf{Aufgabenbereiche}\\
		\midrule
		Project Development & Entwicklung & Projektentwicklung, Site Assessment, Feasibility Studies\\
		\midrule
		Project Management Coordination & Projektleitung & Ansprechpartner gegenüber dem Kunden, Veranlassung der Rechnungsstellung, Projektablaufplanung, Umsetzung der Vertragspflichten, ggf.\ Disposition der Baustelleneinrichtung, ggf.\ Krandisposition\\
		\midrule
		Civil Works & Erdarbeiten & Entwurf, Kontrolle und Abnahme des Fundamentdesigns und des Fundamentbaus; Entwurf, Kontrolle und Abnahme der Zuwegung und der Kranstellfläche\\
		\bottomrule
	\end{tabularx}
	\label{tab:Beispieltab-4_DE}
\end{table}


Tabelle \ref{tab:Beispieltab-5_DE} stellt ein Beispiel für eine mehrseitige Tabelle dar, die mithilfe der 'longtable'-Latex-enviroment erstellt wurde.


\begin{longtable}{lcc}
	\caption{Windenergieproduktion nach Ländern \cite{elec}}\\
	\toprule
	\textbf{Land}             & \textbf{Windenergieerzeugung} & \textbf{Anteil von Gesamtenergie} \\
	& \textbf{(TWh)}                & \textbf{(\%)}                     \\
	\midrule
	\endfirsthead
	\caption{(Fortsetzung) Windenergieproduktion nach Ländern \cite{elec}}\\
	\toprule
	\textbf{Land}             & \textbf{Windenergieerzeugung} & \textbf{Anteil von Gesamtenergie} \\
	& \textbf{(TWh)}                & \textbf{(\%)}                     \\
	\midrule
	\endhead
	World                     & 2098.46                       & 7.3                               \\
	China                     & 762.70                        & 8.6                               \\
	United States             & 434.30                        & 10.1                              \\
	Germany                   & 124.91                        & 22.3                              \\
	United Kingdom            & 80.26                         & 24.6                              \\
	Brazil                    & 80.00                         & 11.8                              \\
	India                     & 70.05                         & 3.8                               \\
	Spain                     & 61.90                         & 21.7                              \\
	France                    & 38.56                         & 8.2                               \\
	Canada                    & 38.13                         & 6.0                               \\
	Turkey                    & 35.14                         & 10.8                              \\
	Sweden                    & 33.33                         & 19.4                              \\
	Australia                 & 31.65                         & 11.6                              \\
	Netherlands               & 21.49                         & 17.7                              \\
	Italy                     & 20.28                         & 7.1                               \\
	Mexico                    & 20.23                         & 5.8                               \\
	Poland                    & 19.73                         & 11.0                              \\
	Denmark                   & 18.93                         & 55.0                              \\
	Norway                    & 15.90                         & 10.4                              \\
	Argentina                 & 14.02                         & 9.8                               \\
	Portugal                  & 13.28                         & 28.3                              \\
	Belgium                   & 12.16                         & 12.8                              \\
	Finland                   & 11.72                         & 16.1                              \\
	Ireland                   & 11.18                         & 33.3                              \\
	Greece                    & 10.87                         & 20.8                              \\
	South Africa              & 9.69                          & 4.5                               \\
	Chile                     & 8.82                          & 10.6                              \\
	Japan                     & 8.17                          & 0.8                               \\
	Vietnam                   & 8.04                          & 3.1                               \\
	Romania                   & 6.97                          & 12.6                              \\
	Austria                   & 6.96                          & 10.6                              \\
	Russia                    & 5.51                          & 0.5                               \\
	Morocco                   & 5.11                          & 12.4                              \\
	Uruguay                   & 4.80                          & 30.5                              \\
	Egypt                     & 4.23                          & 2.0                               \\
	Pakistan                  & 4.12                          & 2.7                               \\
	Taiwan                    & 3.50                          & 1.2                               \\
	South Korea               & 3.37                          & 0.6                               \\
	New Zealand               & 3.20                          & 7.3                               \\
	Thailand                  & 3.00                          & 1.6                               \\
	Kazakhstan                & 2.28                          & 2.0                               \\
	Croatia                   & 2.28                          & 16.1                              \\
	Peru                      & 1.96                          & 3.3                               \\
	Kenya                     & 1.68                          & 13.6                              \\
	Lithuania                 & 1.59                          & 38.0                              \\
	Jordan (2021)             & 1.57                          & 7.2                               \\
	Bulgaria                  & 1.44                          & 2.9                               \\
	Costa Rica                & 1.36                          & 10.8                              \\
	Dominican Republic (2021) & 1.29                          & 7.3                               \\
	Iran                      & 1.25                          & 0.4                               \\
	Philippines               & 0.95                          & 0.8                               \\
	Serbia                    & 0.94                          & 2.7                               \\
	Ukraine                   & 0.83                          & 0.7                               \\
	Honduras (2021)           & 0.80                          & 6.7                               \\
	Panama (2021)             & 0.66                          & 5.9                               \\
	Czech Republic            & 0.65                          & 0.8                               \\
	Estonia                   & 0.63                          & 7.5                               \\
	Nicaragua (2021)          & 0.63                          & 13.6                              \\
	Ethiopia (2021)           & 0.61                          & 4.2                               \\
	Hungary                   & 0.61                          & 1.7                               \\
	Bolivia                   & 0.48                          & 4.3                               \\
	Tunisia                   & 0.42                          & 2.1                               \\
	Mongolia                  & 0.42                          & 5.8                               \\
	Bosnia and Herzegovina    & 0.39                          & 2.4                               \\
	Sri Lanka (2021)          & 0.37                          & 2.3                               \\
	Indonesia                 & 0.36                          & 0.1                               \\
	Guatemala (2021)          & 0.32                          & 2.2                               \\
	Luxembourg                & 0.31                          & 25.2                              \\
	Montenegro                & 0.31                          & 9.4                               \\
	Jamaica (2021)            & 0.30                          & 6.9                               \\
	Senegal (2021)            & 0.25                          & 4.5                               \\
	Cyprus                    & 0.22                          & 4.2                               \\
	Latvia                    & 0.19                          & 3.8                               \\
	\bottomrule
	\label{tab:Beispieltab-5_DE}
\end{longtable}


\clearpage
\chapter{Hinweise zu \LaTeX{}}
\label{cha:latex_DE}

Für die Anfertigung von wissenschaftlichen Arbeiten empfiehlt es sich, auf spezialisierte und etablierte Textsatzsysteme zurückzugreifen. Das IWES empfiehlt \LaTeX{}, da es in der Wissenschaft international weit verbreitet ist. Alle nennenswerten wissenschaftlichen Zeitschriften und Verlage verteilen Formatvorlagen für \LaTeX{}, so dass publizierende Wissenschaftlerinnen und Wissenschaftler nahezu gezwungen sind, sich früher oder später mit diesem System auseinander zu setzen.

\LaTeX{} ist ein kostenfreies OpenSource-Projekt, das sich seit Jahrzehnten etabliert hat. Es ist plattformunabhängig für alle Betriebssysteme (Windows, MacOS, Linux) verfügbar. Im Gegensatz zu Microsoft Word oder verwandten Produkten, in denen innerhalb einer Frontend-Oberfläche gearbeitet wird, handelt es sich bei \LaTeX{} um ein Textsatzsystem, das überwiegend über Backend-Editoren genutzt wird. \LaTeX{} ist eine Auszeichnungssprache. Man bearbeitet Quelldateien in einem Editor. Diese werden dann zu professionellen Textdokumenten kompiliert, die in allen Belangen üblichen Textverarbeitungsprogrammen wie z.\,B. Microsoft Word o.\,Ä. überlegen sind. Insbesondere bei Dokumenten, die mathematische Formeln enthalten, hat sich \LaTeX{} bewährt.

Um mit \LaTeX{} arbeiten zu können, benötigt man drei Zutaten, die im Folgenden für Windows-Computer beschrieben werden (für andere Betriebssysteme kann man geeignete Alternativen über eine Internet-Recherche finden):

\begin{enumerate}
	\item \LaTeX{}-Distribution: Für Windows heißt diese \textit{MiKTeX}. Sie beinhaltet das vollständige Dateisystem, das für das Kompilieren von Dokumenten benötigt wird. Die \LaTeX{}-Distribution sollte zuerst installiert werden, damit der danach installierte Editor die \LaTeX{}-Pfade auf dem Rechner automatisch findet.\\
	\item Editor: Im Editor werden die Quelldateien für das zu kompilierende Dokument bearbeitet. Meist wird ein Backend-Editor genutzt. D.\,h., dass das finale Dokument nicht sofort angezeigt wird, sondern erst kompiliert werden muss. Es gibt zahlreiche Editoren auf dem Markt, die z.\,T. kostenlos verfügbar sind. Es sei beispielhaft der Editor \textit{TeXstudio} genannt, der vielfältige auf \LaTeX{} spezialisierte Funktionen bietet (Syntaxhervorhebung, Autoverfollständigung, etc.) und plattformunabhängig nutzbar ist. \\
	\item Literaturmanagement-Tool: Ein Tool für das Literaturmanagement ist für das wissenschaftliche Arbeiten unumgänglich, unabhängig davon, welche Software letztlich für das Erstellen der wissenschaftlichen Ausarbeitung verwendet wird. Nur mit einem professionellen Tool für das Literaturmanagement ist man langfristig in der Lage, einen Überblick über die Literatur, die man ggf. zitieren möchte, zu behalten. Auch hier gibt es ein breites Angebot an Programmen, die genutzt werden können. \textit{JabRef} ist ein Java-basiertes Tool, das einfach zu nutzen und plattformunabhängig und kostenfrei zur Verfügung steht. Alternativ kann auch auf kostenpflichtige Lösungen mit einem breiteren Funktionsspektrum zurückgegriffen werden, z.\,B. \textit{Citavi}.
\end{enumerate}

Das IWES empfiehlt alternativ zu den oben genannten Programmen das online-Tool \textit{Overleaf}, das die oben genannten Aufzählungspunkte 1. und 2. integriert und vom LUIS als Dienst zur Verfügung gestellt wird. Informationen dazu finden Sie auf der \href{https://www.luis.uni-hannover.de/de/services/speichersysteme/dateiservice/cloud-dienste/overleaf}{\textbf{Website des LUIS}}. Tabelle \ref{tab:programme} fasst die für Windows empfohlenen Programme und ihre Bezugsquellen noch einmal zusammen.

\begin{table}[!htbp]
\normalsize
  \centering
  \caption[Empfohlene \LaTeX{}-Programme für Windows-Systeme.]{Empfohlene \LaTeX{}-Programme für Windows}
    \begin{tabular}{rll}
    	\toprule
    	\textbf{Programm} & \textbf{Funktion} & \textbf{Bezugsquelle} \\
    	\midrule
		Overleaf & Integriertes Online-Tool & \url{https://tex.cloud.uni-hannover.de/}\\
    	MiKTeX & TeX-Distribution für Windows & \url{https://miktex.org/} \\
		TeXstudio & TeX-spezifischer Editor & \url{https://www.texstudio.org/} \\
		JabRef & Literaturmanagement-Tool & \url{https://www.jabref.org/} \\
		\bottomrule
    \end{tabular}%
  \label{tab:programme}%
\end{table}%

Die Quelldateien für diesen Leitfaden sind gleichzeitig die \LaTeX{}-Vorlage für das Erstellen von wissenschaftlichen Arbeiten am IWES. Diese Vorlage soll das wissenschaftliche Arbeiten am IWES vereinfachen und ein einheitliches Layout von studentischen Arbeiten ermöglichen.

Diese \LaTeX{}vorlage ist für die Verwendung des Kompilierungsweges \LaTeX{}$\rightarrow$PDF optimiert. Dies erfordert die Einbindung der Abbildungen und Graphiken als PDF-Vektordateien. Andere Dateiformate werden unter Umständen fehlerhaft oder gar nicht dargestellt. In der Kommandozeile der Editoren muss als Kompilierungsweg daher auch \LaTeX{}$\rightarrow$PDF eingestellt werden. Bei TeXstudio ist dies üblicherweise bereits als Voreinstellung hinterlegt. Bei fehlerhafter Funktionalität sollte dies aber überprüft werden.

Eingebundene Bilder müssen ohne Rand eingebunden werden. Wurde die erstellte PDF-Datei nicht ohne Rand erstellt, so kann dieser mittels Bildbearbeitungsprogramme nachträglich entfernt werden.

Für das Übernehmen von in Excel erstellten Tabellen in \LaTeX{} gibt es ein recht gut funktionierendes Makro namens \textit{Excel2LaTeX}, dass in Excel eingebunden werden kann. Dieses kann unter  \url{http://www.heise.de/download/excel2latex.html} heruntergeladen werden. Alternativ gibt es auch Programme wie \textit{LaTable} oder \textit{Excel-To-LaTeX-Converter}, die das erstellen von Tabellen in \LaTeX{} erleichtern.

\begin{addmargin}[1cm]{1cm}
	\textit{\textbf{Hinweis:} Grafiken und/oder Fotos, die als Abbildungen eingebunden werden sollen, sollten im Format \texttt{<Kapitel-Nr>\_<Abbildungs-Nr>\_<Abbildungs-Name>.<Erweiterung>} benannt werden (Beispiel: \texttt{04\_01\_Beispielbild.pdf} für Abbildung \ref{fig:centered_DE}).}
\end{addmargin}

Bei dem obigen Hinweis gibt die Kapitelnummer das Kapitel an, dem das Bild zugeordnet ist. Sie wird als zweistellige Zahl ausgedrückt (z. B. \texttt{02} für das zweite Kapitel). Die Abbildungsnummer stellt die Reihenfolge der Abbildung im Kapitel dar. Sie wird ebenfalls als zweistellige Zahl ausgedrückt (z.B. \texttt{03} für die dritte Abbildung). Der Abbildungsname beschreibt kurz und knapp und ohne Leerzeichen den Inhalt oder das Thema der Abbildung. Hier ist ein klarer, prägnanter Name, der relevant für den Inhalt der Abbildung ist, erforderlich. Die Dateierweiterung beizeichnet das Dateiformat der Bilddatei. Idealerweise handelt es sich um eine PDF-Datei. In der Regel kann \LaTeX{} aber auch JPG- oder PNG-Dateien verarbeiten.

\cleardoublepage

%=== Verzeichnisse =====================================
\pagenumbering{Roman}
% HINWEIS: Dieser Zähler muss der aktuellen römischen Seitenzahl entsprechen.
% Zählung: Titelseite (I, nicht sichtbar) + Erklärungen + Aufgabenstellung + Kurzfassung + Inhaltsverzeichnis = Seiten vor den Verzeichnissen.
% Wenn Vorspann-Abschnitte hinzugefügt oder entfernt werden, diesen Wert entsprechend anpassen.
\setcounter{page}{7}

% Abbildungs- und Tabellenverzeichnis werden automatisch über listof=totoc
% in der Dokumentklasse zum Inhaltsverzeichnis hinzugefügt — kein \addcontentsline nötig.
\listoffigures
\clearpage

\listoftables
\clearpage

% Symbole und Abkürzungen
% Symbolverzeichnis
\lhead{\slshape{Symbolverzeichnis}}
\addcontentsline{toc}{chapter}{Symbolverzeichnis}

~\\[1.5cm]
{\huge\textbf{Symbolverzeichnis}}
\vspace{1.5cm}

\begin{tabularx}{0.9\textwidth}{l X l}
	\textbf{Symbol} & \textbf{Bedeutung} & \textbf{Einheit} \\
	\hline
	$\alpha$                & Anstellwinkel eines aerodynamischen Profils       & Grad     \\
	$\lambda$              & Schnelllaufzahl          & -        \\
	$\rho_{\text{air}}$    & Luftdichte                                    & kg/m$^3$ \\
	$u_{\text{wind}}$      & Windgeschwindigkeit                           & m/s      \\
	$A$                    & Rotorfläche              & m$^2$    \\
	$c_P$                  & Leistungsbeiwert                         & -        \\
	$c_T$                  & Schubbeiwert                            & -        \\
	$D$                    & Rotordurchmesser                            & m \\
	$P$                    & Leistung        & W        \\
	$R$                    & Rotorradius                                 & m        \\
	$z$                    & Anzahl der Rotorblätter                          & -        \\
\end{tabularx}

\clearpage

% Abkürzungsverzeichnis
\lhead{\slshape{Abkürzungsverzeichnis}}
\addcontentsline{toc}{chapter}{Abkürzungsverzeichnis}

~\\[1.5cm]
{\huge\textbf{Abkürzungsverzeichnis}}
\vspace{1.5cm}

\begin{tabularx}{0.9\textwidth}{l X}
	\textbf{Abkürzung} & \textbf{Bedeutung} \\
	\hline
	AEP & Jährliche Energieproduktion (engl.: Annual Energy Production) \\
	CFD & Numerische Strömungsmechanik (engl.: Computational Fluid Dynamics) \\
	GIS & Geografisches Informationssystem \\
	IWES & Institut für Windenergiesysteme \\
	LUH & Leibniz Universität Hannover \\
	TSR & Schnelllaufzahl (engl.: Tip Speed Ratio) \\
	WEA & Windenergieanlage \\
\end{tabularx}

\clearpage


%=== Literaturverzeichnis ==============================
\sloppy
\addtocontents{toc}{\vspace{8ex}}
\printbibliography[heading=bibintoc, title=Literaturverzeichnis]
\cleardoublepage

%=== Anhang ============================================
\appendix
% Anhang
\chapter{Erster Anhang}
Dies ist der erste Anhang.

\chapter{Zweiter Anhang}
Dies ist der zweite Anhang.

\clearpage


\end{document}
